\documentclass[11pt,openany]{book}

\usepackage[
  paperwidth=6.5in,
  paperheight=9.5in,
  inner=0.82in,
  outer=0.68in,
  top=0.78in,
  bottom=0.84in,
  headheight=25pt,
  headsep=16pt,
  footskip=23pt
]{geometry}
\usepackage{fontspec}
\defaultfontfeatures{Ligatures=TeX}
\setmainfont[
  Path=solutions_ja/fonts/,
  Extension=.ttf,
  UprightFont=*-Regular,
  BoldFont=*-Bold,
  ItalicFont=*-Regular,
  BoldItalicFont=*-Bold,
  ItalicFeatures={FakeSlant=0.18},
  BoldItalicFeatures={FakeSlant=0.18}
]{NotoSerifJP}
\XeTeXlinebreaklocale "ja"
\XeTeXlinebreakskip=0pt plus 1pt
\XeTeXlinebreakpenalty=0

\usepackage[nocomments,norevision]{omphys}
\usepackage{amsthm}
\usepackage{booktabs,longtable,tabularx}
\usepackage{enumitem}
\usepackage[most]{tcolorbox}
\usepackage{fancyhdr}
\usepackage{titlesec}
\usepackage{hyperref}
\hypersetup{
  unicode=true,
  pdftitle={量子共鳴を大域接続問題として理解する：完全解答集},
  pdfauthor={森川奥統・小川将也},
  colorlinks=true,
  linkcolor=blue!45!black,
  urlcolor=blue!55!black
}
\setlength{\emergencystretch}{2em}
\allowdisplaybreaks[2]

\pagestyle{fancy}
\fancyhf{}
\fancyhead[LE]{\small\thepage\quad\nouppercase{\leftmark}}
\fancyhead[RO]{\small\nouppercase{\rightmark}\quad\thepage}
\renewcommand{\headrulewidth}{0.35pt}
\fancypagestyle{plain}{
  \fancyhf{}
  \fancyfoot[C]{\thepage}
  \renewcommand{\headrulewidth}{0pt}}
\titleformat{\chapter}[display]
  {\normalfont\huge\bfseries}
  {\color{teal!55!black}第\thechapter 部}
  {0.55ex}{}
\titlespacing*{\chapter}{0pt}{-12pt}{22pt}

\newtcolorbox{problemstatement}{
  enhanced,breakable,colback=teal!3,colframe=teal!55!black,
  title=問題,fonttitle=\bfseries}
\newtcolorbox{solutionbox}{
  enhanced,breakable,colback=blue!2,colframe=blue!48!black,
  title=解答,fonttitle=\bfseries}
\newtcolorbox{checkbox}{
  enhanced,breakable,colback=orange!5,colframe=orange!70!black,
  title=検算・数値確認,fonttitle=\bfseries}
\newcommand{\solproblem}[3]{%
  \section*{#1}%
  \phantomsection%
  \addcontentsline{toc}{section}{#1}%
  \markright{#1}%
  \begin{problemstatement}#2\end{problemstatement}%
  \begin{solutionbox}#3\end{solutionbox}%
}
\newcommand{\oplusint}{\int^{\oplus}}

\title{量子共鳴を大域接続問題として理解する\\[0.8ex]
  \Large 問題・演習 完全解答集}
\author{森川億人 \and 小川翔也}
\date{2026年8月}

\begin{document}
\frontmatter
\maketitle
\chapter*{本解答集の使い方}
本冊子は、英語版教科書 \textit{Quantum Resonance as a Global
Connection Problem} の Stage I--VIII の99問と、再総和付録の7問を
すべて解く。問題番号は本体と一致する。各解答では、結論だけでなく、
符号、作用素の定義域、解析接続するシート、境界条件、数値計算の収束判定
まで明記する。数値実験を要求する問題については、再現可能なアルゴリズムと
合否基準を解答の一部とした。

式の記法は本体に合わせる。特に
$\Resolvent(z)=(z-H)^{-1}$、虚数単位 $\rmi$、微分記号 $\rmd$、
積分測度 $\dd x$、入射接続係数 $\Cincoming$ を用いる。
\tableofcontents

\mainmatter
\chapter{Stage I：作用素・直積分・Fredholm 理論}

\solproblem{I.1　effect と POVM}{
$0\leq F\leq\identity$ に対する Born 則が $[0,1]$ の確率を与えることを示し、射影値でない二値 POVM を作れ。}{
密度作用素を $\rho\geq0$, $\Tr\rho=1$ とする。正作用素の順序は
$0\leq \rho^{1/2}F\rho^{1/2}\leq\rho$ を与えるから、
\begin{equation}
  0\leq p(F)=\Tr(\rho F)
  =\Tr(\rho^{1/2}F\rho^{1/2})\leq\Tr\rho=1.
\end{equation}
二準位系で $0<\eta<1$ とし
\begin{equation}
  F_1=\eta\ket{1}\bra{1},\qquad F_0=\identity-F_1
\end{equation}
と置けば、$F_0+F_1=\identity$ だが $F_1^2=\eta F_1\neq F_1$ なので
PVM ではない。励起状態を効率 $\eta$ でのみ検出する有限効率検出器である。
より対称な有限分解能測定は、鋭い射影 $P$ を
$F_\pm=(\identity\pm\eta(2P-\identity))/2$ へぼかして得られる。}

\solproblem{I.2　区間上の運動量の自己共役拡張}{
$P=-\rmi\hbar\rmd/\rmd x$ on $(0,L)$ に境界条件
$\psi(L)=\rme^{\rmi\alpha}\psi(0)$ を課し、自己共役性と固有値を求めよ。}{
最大作用素の関数 $\phi,\psi\in H^1(0,L)$ に対して
\begin{equation}
 \dualpair{\phi}{P\psi}-\dualpair{P\phi}{\psi}
 =-\rmi\hbar[\phi^*(L)\psi(L)-\phi^*(0)\psi(0)].
\end{equation}
両者が同じ準周期条件を満たせば右辺は0である。逆に、$P^*$ の定義域に
属する $\phi$ について右辺が全ての $\psi$ に対して0になるには
$\phi(L)=\rme^{\rmi\alpha}\phi(0)$ が必要である。したがって
$D(P^*)=D(P)$ で自己共役である。$P\psi=p\psi$ の解は
$\psi=C\rme^{\rmi px/\hbar}$ であり、
\begin{equation}
  \rme^{\rmi pL/\hbar}=\rme^{\rmi\alpha},\qquad
  p_n=\frac{\hbar}{L}(2\pi n+\alpha),\quad n\in\bbZ,
\end{equation}
$\psi_n=L^{-1/2}\rme^{\rmi p_nx/\hbar}$ となる。境界位相が作用素の
定義域、したがってスペクトルの一部であることが要点である。}

\solproblem{I.3　弱収束と確率の逃散}{
正規直交列 $e_n$ が0へ弱収束するがノルム収束しないことを示せ。}{
任意の $f\in\Hil$ に対し Bessel の不等式
$\sum_n|\dualpair{f}{e_n}|^2\leq\|f\|^2$ が成り立つので、各項は0へ
収束する。従って $e_n\rightharpoonup0$。一方 $\|e_n\|=1$ だから
ノルム収束しない。$e_n(x)=\phi(x-n)$ かつ $\phi\in L^2(\real)$ なら、
任意の固定コンパクト集合 $K$ について
\begin{equation}
 \int_K|e_n(x)|^2\dd x=\int_{K-n}|\phi(y)|^2\dd y\longrightarrow0.
\end{equation}
全確率は1のままだが、有限位置の検出器から確率が逃げる。}

\solproblem{I.4　運動量表示からエネルギー直積分へ}{
$\rmd^3p=p^2\dd p\rmd\Omega$ と $E=p^2/(2\mu)$ からユニタリな
エネルギー表示を導け。}{
$\dd E=(p/\mu)\dd p$、従って $p^2\dd p=\mu p\dd E$ である。
$\psi(\bm p)=\psi(p,\Omega)$ に対し
\begin{equation}
 (U\psi)(E,\Omega)=(\mu p)^{1/2}\psi(p,\Omega),
 \qquad p=\sqrt{2\mu E}
\end{equation}
と定めると
\begin{align}
 \|U\psi\|^2
 &=\int_0^\infty\dd E\int\rmd\Omega\,
     \mu p|\psi(p,\Omega)|^2 \\
 &=\int_0^\infty p^2\dd p\int\rmd\Omega\,
     |\psi(p,\Omega)|^2=\|\psi\|^2.
\end{align}
従って $U$ は $L^2(\real^3,\rmd^3p)$ から
$\oplusint_{[0,\infty)}L^2(S^2,\rmd\Omega)\dd E$ へのユニタリ写像で、
$H_0$ は各 fiber 上で $E$ の乗算になる。}

\solproblem{I.5　multiplicity-one の commutant}{
全ての乗算作用素と可換な有界作用素を決定せよ。}{
$A M_f=M_fA$ とする。特に任意の可測集合 $\Delta$ の
$M_{\chi_\Delta}$ と可換なので、$A$ は $\Delta$ に台を持つ関数を同じ
$\Delta$ 内へ写す。有限測度の場合 $g=A\identity$ と置くと、単関数
$s=\sum_jc_j\chi_{\Delta_j}$ に対し
\begin{equation}
 As=\sum_jc_jM_{\chi_{\Delta_j}}A\identity=gs.
\end{equation}
単関数の稠密性から $A=M_g$、また特性関数で評価して
$g\in L^\infty$ かつ $\|g\|_\infty\leq\|A\|$。sigma 有限の場合は有限
測度部分に制限して貼り合わせる。fiber 次元が $m(E)>1$ なら、同じ
議論がエネルギーを混ぜないことだけを強制し、
$A=\oplusint A(E)\dd\mu(E)$、$A(E)\in\Banach(\complex^{m(E)})$ となる。}

\solproblem{I.6　スペクトル測度の二例}{
二準位 Hamiltonian と一次元自由粒子 wave packet のスペクトル測度を求めよ。}{
$H=E_1P_1+E_2P_2$、$\psi=c_1\ket{1}+c_2\ket{2}$ なら
\begin{equation}
 \mu_\psi(\Delta)=\dualpair{\psi}{P_H(\Delta)\psi}
 =|c_1|^2\chi_\Delta(E_1)+|c_2|^2\chi_\Delta(E_2),
\end{equation}
すなわち $\rmd\mu_\psi=|c_1|^2\delta(E-E_1)\dd E+
|c_2|^2\delta(E-E_2)\dd E$。
一次元では $E=p^2/(2m)$ に $p=\pm\sqrt{2mE}$ の二枝があり、
\begin{equation}
 \frac{\rmd\mu_\psi}{\rmd E}
 =\frac{m}{\sqrt{2mE}}
 \left(|\psi(\sqrt{2mE})|^2+|\psi(-\sqrt{2mE})|^2\right).
\end{equation}
二枝がエネルギー fiber の multiplicity 2 を表す。}

\solproblem{I.7　resolvent の解析性}{
resolvent 恒等式から $\Resolvent(z)$ の作用素ノルム解析性を示せ。}{
直接乗算して
\begin{equation}
 \Resolvent(z)-\Resolvent(w)=(w-z)\Resolvent(z)\Resolvent(w)
\end{equation}
を得る。$|h|\|\Resolvent(z)\|<1$ なら
\begin{equation}
 \Resolvent(z+h)=\Resolvent(z)
 [\identity+h\Resolvent(z)]^{-1}
 =\Resolvent(z)-h\Resolvent(z)^2+\order(h^2),
\end{equation}
従って $\partial_z\Resolvent(z)=-\Resolvent(z)^2$。ただし連続スペクトル
上では $z-H$ が有界逆を持たず、この局所級数の円板は切断を越えない。
解析接続には重み付き空間、解析 vector、Jost 表現など追加構造が要る。}

\solproblem{I.8　Feshbach 作用素と微分 metric}{
$P+Q=\identity$ を用いて有効 Hamiltonian と極の規格化因子を導け。}{
$(E-H)\Psi=0$ を $P,Q$ で射影すると
\begin{align}
 (E-H_{PP})P\Psi-H_{PQ}Q\Psi&=0,\\
 -H_{QP}P\Psi+(E-H_{QQ})Q\Psi&=0.
\end{align}
$E-H_{QQ}$ が可逆なシートでは
$Q\Psi=(E-H_{QQ})^{-1}H_{QP}P\Psi$。従って
\begin{equation}
 [E-\vsub{H}{eff}(E)]P\Psi=0,
 \quad \vsub{H}{eff}(E)=H_{PP}+H_{PQ}(E-H_{QQ})^{-1}H_{QP}.
\end{equation}
$D(E)=E-\vsub{H}{eff}(E)$ の単純零点 $\resonantE$ 近傍で、右・左零 vector
$\ket{r},\bra{l}$ を用いると
\begin{equation}
 D(E)^{-1}\sim
 \frac{\ket{r}\bra{l}}
 {(E-\resonantE)\bra{l}[\identity-\partial_E\vsub{H}{eff}(\resonantE)]\ket{r}}.
\end{equation}
分母の微分 metric は、消去した $Q$ 空間成分のエネルギー依存を含めた
全状態の規格化である。}

\solproblem{I.9　閾値付き Lorentzian の生存振幅}{
下端閾値を持つ Lorentzian 密度の Fourier 積分を極項と閾値項に分けよ。}{
例として
\begin{align}
 \rho(E)&=\Theta(E-\vsub{E}{th})
 \frac{\Gamma/(2\pi)}{(E-E_0)^2+(\Gamma/2)^2},\\
 A(t)&=\int_{\vsub{E}{th}}^\infty
 \rho(E)\rme^{-\rmi Et/\hbar}\dd E.
\end{align}
を取る。$t>0$ で下半平面へ閉じると
$\resonantE=E_0-\rmi\Gamma/2$ の留数が
$Z\rme^{-\rmi\resonantE t/\hbar}$ を与える。閾値から延びる輪郭は端点
積分であり、部分積分により
\begin{equation}
 \vsub{A}{th}(t)=
 \frac{\hbar\rho(\vsub{E}{th})}{\rmi t}
 \rme^{-\rmi \vsub{E}{th}t/\hbar}+\order(t^{-2}).
\end{equation}
密度が $(E-\vsub{E}{th})^\alpha$ で始まれば
$\vsub{A}{th}\sim t^{-\alpha-1}\rme^{-\rmi \vsub{E}{th}t/\hbar}$。従って
中間時間は指数減衰、十分長時間は閾値由来の冪減衰となる。}

\solproblem{I.10　四つの収束位相}{
ノルム、強作用素、弱作用素、強 resolvent 収束の例を一つずつ挙げよ。}{
次が典型例である。
\begin{enumerate}[label=(\alph*)]
 \item $A_n=\identity/n\to0$ は作用素ノルム収束。
 \item $\ell^2$ の最初の $n$ 成分への射影 $P_n\to\identity$ は強収束
  だが $\|\identity-P_n\|=1$。
 \item 片側 shift $S^n\to0$ は弱作用素収束だが
  $\|S^ne_1\|=1$ なので強収束しない。
 \item $L^2(\real)$ 上の乗算作用素
  $H_n(x)=\max(-n,\min(x,n))$ は $H_n$ 自身について一様有界ではないが、
  任意の $z\notin\real$ で
  $(z-H_n)^{-1}\to(z-x)^{-1}$ が優収束定理により強収束する。
\end{enumerate}
強収束をノルム収束と誤れば有限 detector での収束から全状態一様誤差を
主張してしまう。弱収束を強収束と誤れば逃散する確率を消失とみなし、強
resolvent 収束を作用素収束と誤れば非有界 observable の domain 情報を失う。}

\solproblem{I.11　compact 作用素への強収束の一様化}{
$P_N\to\identity$ strongly、$K$ compact なら
$\|(\identity-P_N)K\|\to0$ を示せ。}{
$K$ が単位球を写した集合の閉包は compact なので、任意の $\varepsilon>0$
に対し有限個の点 $y_1,\dots,y_m$ の $\varepsilon$ 球で覆える。$P_N$ は
一様有界であり、有限個の $y_j$ について
$\max_j\|(\identity-P_N)y_j\|\to0$。任意の $\|x\|\leq1$ に対し
$\|Kx-y_j\|<\varepsilon$ なる $j$ を選べば
\begin{equation}
 \|(\identity-P_N)Kx\|
 \leq\|\identity-P_N\|\varepsilon
 +\max_j\|(\identity-P_N)y_j\|.
\end{equation}
上限を取り $N\to\infty$、次に $\varepsilon\to0$ とする。$K=\identity$
なら $\|\identity-P_N\|=1$ で反例になる。}

\solproblem{I.12　自由 Green kernel の Hilbert--Schmidt 性}{
三次元の off-shell 自由 Green kernel と有界 compact support potential から
作る Birman--Schwinger kernel の二乗可積分性を示せ。}{
$z=k^2/(2\mu)$、$\im k>0$ とし、定数を除けば
\begin{equation}
 G_0(z;\bm x,\bm y)=
 \frac{\rme^{\rmi k|\bm x-\bm y|}}{|\bm x-\bm y|},\quad
 K=|V|^{1/2}G_0\operatorname{sgn}(V)|V|^{1/2}.
\end{equation}
対角 $r=|\bm x-\bm y|\downarrow0$ では
$|G_0|^2\sim r^{-2}$ だが体積要素は $4\pi r^2\dd r$ なので局所積分は
$\int_0^\epsilon\dd r<\infty$。$V$ の台が compact なら両変数は有界領域
に制限され、残りは連続で有限である。台を緩めても $\im k>0$ なら
$\rme^{-2\im k r}$ が大距離を抑える。従って $K$ は
Hilbert--Schmidt、特に compact である。}

\solproblem{I.13　Birman--Schwinger の符号}{
本書の resolvent 規約で Birman--Schwinger 同値を証明し、逆規約との符号差を示せ。}{
$(H_0+V-z)\psi=0$ は
$(z-H_0)\psi=V\psi$、従って
$\psi=\Resolvent_0(z)V\psi$ と同値である。極分解
$V=|V|^{1/2}J|V|^{1/2}$ を用い
$\phi=|V|^{1/2}\psi$ と置けば
\begin{equation}
 \phi=K(z)\phi,\qquad
 K(z)=|V|^{1/2}\Resolvent_0(z)J|V|^{1/2}.
\end{equation}
逆向きは $\psi=\Resolvent_0J|V|^{1/2}\phi$ と構成する。
$\widetilde R_0=(H_0-z)^{-1}=-\Resolvent_0$ を採用すれば
$\phi=-\widetilde K(z)\phi$、すなわち固有値条件は
$-1\in\spectrum(\widetilde K)$ になる。}

\solproblem{I.14　解析 Fredholm 逆の留数}{
$A(z)=\identity-K(z)$ の単純零点における $A(z)^{-1}$ の Laurent 係数を求めよ。}{
$A(z_0)\ket{r}=0$、$\bra{l}A(z_0)=0$、
$d=\bra{l}A'(z_0)\ket{r}\neq0$ とする。Ansatz
$A^{-1}=(z-z_0)^{-1}B_{-1}+B_0+\cdots$ を $AA^{-1}=\identity$ に代入すると
$A(z_0)B_{-1}=0$ なので $B_{-1}=c\ket{r}\bra{l}$。定数項を左から
$\bra{l}$ で挟むと $c d\bra{l}=\bra{l}$、従って
\begin{equation}
 A(z)^{-1}=\frac{\ket{r}\bra{l}}
 {(z-z_0)\bra{l}A'(z_0)\ket{r}}+\order(1).
\end{equation}
$\ket{r}\mapsto c\ket{r}$、$\bra{l}\mapsto c^{-1}\bra{l}$ では分子・分母
とも不変である。}

\solproblem{I.15　乗算作用素は compact でない}{
$L^2([0,1],\dd E)$ 上の非零乗算作用素が compact でないことを示せ。}{
$f\neq0$ a.e. でないなら、ある $\varepsilon>0$ に対し
$A=\{|f|\geq\varepsilon\}$ が正測度を持つ。$A$ を互いに素な正測度集合
$A_n$ に分け、$e_n=\chi_{A_n}/\sqrt{\mu(A_n)}$ とする。$e_n$ は正規直交し、
$M_fe_n$ の台も互いに素で $\|M_fe_n\|\geq\varepsilon$。従って像列は
収束部分列を持たず、$M_f$ は compact でない。これに対し局在
Birman--Schwinger 作用素は積分 kernel により変数を混ぜて平滑化し、
対角特異性も二乗可積分である。エネルギーごとの乗算と、座標間を結ぶ
compact kernel は異なる作用素構造であり矛盾しない。}

\chapter{Stage II：散乱、Jost 関数、長距離相関}

\solproblem{II.1　二体問題の分離}{
二粒子の重心・相対運動を分離し、domain が tensor 積に分かれる仮定を述べよ。}{
$M=m_1+m_2$、$\mu=m_1m_2/M$、
$\bm R=(m_1\bm r_1+m_2\bm r_2)/M$、$\bm r=\bm r_1-\bm r_2$ と置く。
共役運動量を $\bm P=\bm p_1+\bm p_2$、
$\bm p=(m_2\bm p_1-m_1\bm p_2)/M$ とすれば正準変換であり、
\begin{equation}
 \frac{\bm p_1^2}{2m_1}+\frac{\bm p_2^2}{2m_2}
 =\frac{\bm P^2}{2M}+\frac{\bm p^2}{2\mu}.
\end{equation}
$V=V(\bm r)$ が並進不変で、相対 kinetic energy に関して相対 form 有界
（相対 bound $<1$ で十分）なら内部 Hamiltonian
$h=\bm p^2/(2\mu)+V$ は自己共役となり、
$H=\vsub{H}{cm}\otimes\identity+\identity\otimes h$ を tensor 積 domain
（厳密には閉 quadratic form の和）上に定義できる。外場、速度依存相互作用、
重心に依存する境界条件があればこの分離は成立しない。}

\solproblem{II.2　Lippmann--Schwinger 方程式の符号}{
本書の resolvent と時間規約から outgoing Lippmann--Schwinger 方程式を導け。}{
$(H_0+V)\ket{\psi^{(+)}}=E\ket{\psi^{(+)}}$ より
$(E-H_0)\ket{\psi^{(+)}}=V\ket{\psi^{(+)}}$。自由解を加え、limiting
absorption を取ると
\begin{equation}
 \ket{\psi^{(+)}}=\ket{\phi}
 +\Resolvent_0(E+\rmi0)V\ket{\psi^{(+)}}.
\end{equation}
時間依存を $\rme^{-\rmi Et/\hbar}$ としたので、$+\rmi0$ の kernel は
\begin{equation}
 \Resolvent_0(E+\rmi0;\bm r,\bm r')
 =-\frac{\mu}{2\pi\hbar^2}
   \frac{\rme^{\rmi k|\bm r-\bm r'|}}{|\bm r-\bm r'|},
\end{equation}
すなわち outgoing wave である。$(H_0-E-\rmi0)^{-1}$ を使う文献では
resolvent 自体に負号が入る。}

\solproblem{II.3　三次元自由 Green 関数と Born 振幅}{
自由 Green kernel を評価し、第一 Born 振幅を求めよ。}{
Fourier 積分を角度積分し、$p$ 平面の pole を放射条件に従って閉じると
\begin{equation}
 G_0^{(+)}(\bm r-\bm r')
 =-\frac{\mu}{2\pi\hbar^2}
 \frac{\rme^{\rmi k|\bm r-\bm r'|}}{|\bm r-\bm r'|}.
\end{equation}
$r\to\infty$ では
$|\bm r-\bm r'|=r-\widehat{\bm r}\cdot\bm r'+\order(r^{-1})$。従って
\begin{equation}
 \psi^{(+)}(\bm r)\sim \rme^{\rmi\bm k_i\cdot\bm r}
 +\frac{\rme^{\rmi k_fr}}{r}f(\bm k_f,\bm k_i),
\end{equation}
第一反復で
\begin{equation}
 \vsub{f}{B}(\bm k_f,\bm k_i)
 =-\frac{\mu}{2\pi\hbar^2}
 \int\rme^{-\rmi\bm k_f\cdot\bm r}
 V(\bm r)\rme^{\rmi\bm k_i\cdot\bm r}\dd^3r.
\end{equation}}

\solproblem{II.4　多 channel 断面積の flux}{
確率流と $k_f/k_i$ factor を導き、flux 規格化した $S$ 行列を説明せよ。}{
Schrödinger 方程式とその複素共役の差から
\begin{equation}
 \bm j=\frac{\hbar}{2\mu\rmi}
 (\psi^*\bm\nabla\psi-\psi\bm\nabla\psi^*)
\end{equation}
を得る。channel $c$ の球面波 $A_c\rme^{\rmi k_cr}/r$ の radial flux は
$j_r=\hbar k_c|A_c|^2/(\mu_cr^2)$。従って同じ平面波振幅規格化なら
\begin{equation}
 \frac{\rmd\sigma_{fi}}{\rmd\Omega}
 =\frac{v_f}{v_i}|f_{fi}|^2
 =\frac{k_f\mu_i}{k_i\mu_f}|f_{fi}|^2,
\end{equation}
同じ reduced mass なら $k_f/k_i$。一方、基底を
$\vsub{\ket{c,E}}{flux}=v_c^{-1/2}\ket{c,k_c}$ と規格化すれば速度因子は
振幅に吸収され、on-shell $S(E)^dagger S(E)=\identity$ は通常の行列積で
書ける。}

\solproblem{II.5　光学定理と吸収}{
partial-wave unitarity から光学定理を示し、複素 optical potential の場合を論ぜよ。}{
弾性一 channel で $S_l=\rme^{2\rmi\delta_l}$ とすると
\begin{equation}
 f(\theta)=\frac{1}{2\rmi k}
 \sum_l(2l+1)(S_l-1)P_l(\cos\theta).
\end{equation}
直交性から
$\vsub{\sigma}{el}=(\pi/k^2)\sum_l(2l+1)|S_l-1|^2$。$|S_l|=1$ を使えば
\begin{equation}
 \vsub{\sigma}{tot}=\vsub{\sigma}{el}
 =\frac{4\pi}{k}\im f(0).
\end{equation}
$V=V_R+\rmi W$、$W\leq0$ では $|S_l|\leq1$ となり、欠損 flux は
\begin{equation}
 \vsub{\sigma}{reac}=\frac{\pi}{k^2}
 \sum_l(2l+1)(1-|S_l|^2)
 =-\frac{2}{\hbar v}\int W(\bm r)|\psi^{(+)}|^2\dd^3r\geq0.
\end{equation}
光学定理の全断面積は $\vsub{\sigma}{el}+\vsub{\sigma}{reac}$ である。}

\solproblem{II.6　Jost 解による radial Green 関数}{
regular 解と outgoing Jost 解から Green 関数を作り、単純零点の留数を求めよ。}{
reduced radial 方程式を $[E-h_l]u=0$ とし、origin で regular な
$\varphi_l(k,r)$ と infinity で outgoing な $f_l^{(+)}(k,r)$ を取る。
Wronskian $W_l(k)=W[\varphi_l,f_l^{(+)}]$ は $r$ に依らない。微分 jump
を合わせれば
\begin{equation}
 G_l^{(+)}(r,r';k)=\frac{2\mu}{\hbar^2}
 \frac{\varphi_l(k,r_<)f_l^{(+)}(k,r_>)}{W_l(k)}.
\end{equation}
規約により Jost 関数は $W_l$ の非零定数倍である。$W_l(k_R)=0$、
$W_l'(k_R)\neq0$ なら両解は同じ Gamow 解 $u_R$ に比例し、
\begin{equation}
 \Res_{k=k_R}G_l^{(+)}(r,r';k)
 =\frac{2\mu}{\hbar^2}
 \frac{u_R(r)u_R(r')}{W_l'(k_R)}
\end{equation}
（比例定数は採用した Jost 規格化へ吸収）となる。}

\solproblem{II.7　Coulomb 対数位相}{
Coulomb potential を漸近直線軌道に沿って積分し、plane wave が不適切な理由を示せ。}{
速度 $v$、impact parameter $b$ の直線 $r(t)=\sqrt{b^2+v^2t^2}$ に沿う
eikonal 位相は
\begin{equation}
 \chi(T)=-\frac{1}{\hbar}\int^T\frac{g}{r(t)}\dd t
 =-\frac{g}{\hbar v}\operatorname{arsinh}\frac{vT}{b}
 =-\eta\log\frac{2vT}{b}+o(1).
\end{equation}
位相は定数へ収束せず、radial Coulomb 波は
$\rme^{\rmi[kr-\eta\log(2kr)+\sigma_l]}$ を含む。従って plane wave を
基準に差を有限定数として定義する通常の散乱振幅は存在しない。Coulomb
波を comparison state に取った振幅と plane-wave Born 振幅は同じ量では
ない。}

\solproblem{II.8　Dollard modifier}{
通常の Møller 極限が Coulomb 位相で失敗することを示し、modifier を構成せよ。}{
free trajectory 上の相互作用表示の積分は
$\int^t g/(v|s|)\dd s=(g/v)\operatorname{sgn}(t)\log|t|+O(1)$。
従って $\rme^{\rmi Ht/\hbar}\rme^{-\rmi H_0t/\hbar}$ は対数的に回転し続け、
強極限を持たない。運動量表示で $v(p)=p/\mu$ とし
\begin{equation}
 U_D(t)=\rme^{-\rmi H_0t/\hbar}
 \rme^{-\frac{\rmi}{\hbar}\frac{g}{v(p)}
       \operatorname{sgn}(t)\log|t|}
\end{equation}
を comparison dynamics とする。修正 wave operator
$\Omega_D^{(\pm)}=\operatorname{s-lim}_{t\to\mp\infty}
\rme^{\rmi Ht/\hbar}U_D(t)$ では発散位相が相殺される。有限定数や
reference time の違いは asymptotic state の位相規約である。}

\solproblem{II.9　三荷電粒子の漸近領域}{
三つの cluster 漸近領域を列挙し、pairwise Coulomb 位相が一様でない理由を述べよ。}{
Jacobi 座標 $(\bm x_\alpha,\bm y_\alpha)$ に対し、(i) 粒子1と2が有限距離で
cluster を作り粒子3が遠い領域、(ii) 2と3が cluster、(iii) 3と1が
cluster、に加え三距離が同時に大きい genuine three-body breakup 領域が
ある。各 cluster 領域では大変数は spectator 座標 $y_\alpha$、内部座標
$x_\alpha$ は有限である。一方 breakup では hyperradius
$\rho=(x_\alpha^2+y_\alpha^2)^{1/2}$ が大きく hyperangle も動く。
各 pair の $\eta_{ij}\log r_{ij}$ を単純に足した形は、cluster 境界で
$r_{ij}$ の大小が入れ替わると誤差の次数が一様でなく、overlap 領域の
相関振幅と三体位相を決めない。これが一つの elementary ``Coulomb plane
wave'' が知られていない核心である。}

\solproblem{II.10　soft-photon coherent dressing}{
coherent profile $f(\bm k)$ が Fock vector を与える条件を求め、Coulomb field と比較せよ。}{
coherent state
$\ket{f}=\rme^{-\|f\|^2/2}\rme^{a^\dagger(f)}\ket{0}$ が Fock 空間に
属する必要十分条件は
\begin{equation}
 \|f\|^2=\sum_\lambda\int|f_\lambda(\bm k)|^2\dd^3k<\infty.
\end{equation}
charged particle の long-range field を再現する dressing は soft limit で
$f(\bm k)\sim |\bm k|^{-3/2}$（角度因子を除く）を持つので
$\|f\|^2\sim\int_0^\Lambda k^2k^{-3}\dd k=int_0^\Lambda\dd k/k$ と
対数発散する。従って物理的 charged sector は真空 Fock 表現と unitary
同値でなく、有限 soft cutoff を外す順序を指定しなければならない。}

\solproblem{II.11　Lagrange 恒等式}{
一般 Sturm--Liouville 表式の Lagrange 恒等式と endpoint bracket を示せ。}{
$\tau u=w^{-1}[-(pu')'+qu]$ とし
$[u,v]=p(u^*v'-u'^*v)$ と定める。積の微分を展開すると
\begin{equation}
 \int_a^b\{u^*(\tau v)-(\tau u)^*v\}w\dd x
 =[u,v](b)-[u,v](a).
\end{equation}
regular endpoint $a$ の separated 条件を
$\cos\alpha\,u(a)+\sin\alpha\,p(a)u'(a)=0$ と書く。同じ $\alpha$ を
満たす二関数の境界 data は同じ一次元部分空間に属するため、その
symplectic determinant $[u,v](a)$ は0である。}

\solproblem{II.12　逆二乗 potential の limit-point 閾値}{
$-u''+g u/r^2=0$ の Frobenius 指数を求めよ。}{
$u\sim r^\nu$ を代入すると $\nu(\nu-1)=g$、従って
\begin{equation}
 \nu_\pm=\frac{1\pm\sqrt{1+4g}}{2}.
\end{equation}
$\int_0^\epsilon r^{2\nu}\dd r<\infty$ は $\nu>-1/2$ と同値。
$\nu_+$ は常に可積分で、$\nu_-$ も可積分なのは
$\sqrt{1+4g}<2$、すなわち $g<3/4$。従って $g\geq3/4$ では origin は
limit point、$g<3/4$ では（下方有界性 $g>-1/4$ の範囲で）limit circle
で境界条件が必要。三次元 physical $s$ wave は $g=0$ だが元の波動関数
$R=u/r$ の正則性を課し Friedrichs 条件 $u(0)=0$ を選ぶ。}

\solproblem{II.13　Sturm--Liouville Green kernel の jump}{
Green kernel の符号と規格化を jump 条件から決めよ。}{
$u_a$ を左境界条件、$u_b$ を右境界条件を満たす解とし、
$W_p=p(u_a u_b'-u_a'u_b)$ とする。$(z-L)G=\delta/w$ の規約では
\begin{equation}
 G(x,x';z)=-\frac{u_a(x_<)u_b(x_>)}{W_p(z)}
\end{equation}
（$L=w^{-1}[-(p\partial_x)'+q]$）となる。実際、$x=x'$ をまたいで
積分すると $p[G'_+-G'_-]=-1$ であり、上式の jump がこれに一致する。
$(L-z)^{-1}$ を resolvent とすれば右辺の delta と jump の符号が逆転し、
kernel 全体に負号が付く。}

\solproblem{II.14　球形井戸の Jost 関数}{
matching により Jost 関数と resonance 条件を求め、bound-state zero を閾値まで追跡せよ。}{
$U=-U_0$ for $r<a$、$U=0$ for $r>a$ とし
$q=(k^2+U_0)^{1/2}$。regular 解は $\varphi=\sin(qr)/q$。外側を
$A\rme^{\rmi kr}+B\rme^{-\rmi kr}$ として $r=a$ で値と微分を合わせると
\begin{align}
 2B\rme^{-\rmi ka}&=\frac{\sin qa}{q}+\frac{\rmi}{k}\cos qa,\\
 2A\rme^{\rmi ka}&=\frac{\sin qa}{q}-\frac{\rmi}{k}\cos qa.
\end{align}
従って一つの規格化は
\begin{equation}
 F(k)=\rme^{\rmi ka}\left(\cos qa-
 \rmi\frac{k}{q}\sin qa\right),
 \qquad B=0\Longleftrightarrow q\cot qa=\rmi k.
\end{equation}
bound state は $k=\rmi\kappa$、$\kappa>0$ なので
$q\cot qa=-\kappa$。$U_0$ を下げると最上位の $\kappa$ は0へ近づき、
$k=0$ で scattering length が発散する。その後零点は負虚軸の非物理
sheet へ移り virtual state となる。数値追跡では $q$ の branch を連続に
保つ。}

\solproblem{II.15　Wronskian 微分と Gamow 規格化}{
radial 方程式を $k$ で微分して積分恒等式を導け。}{
$[-\partial_r^2+U(r)-k^2]u(k,r)=0$ を $k$ で微分すると
\begin{equation}
 [-\partial_r^2+U-k^2]\partial_ku=2ku.
\end{equation}
元の式を $\partial_ku$ 倍したものとの差を $[0,R]$ で積分して
\begin{equation}
 2k\int_0^Ru(k,r)^2\dd r
 =W[u,\partial_ku](R)-W[u,\partial_ku](0)
\end{equation}
を得る。複素共役は取らない。regular 規格化なら origin 項は既知で、
bound state では infinity の減衰により $R\to\infty$ の境界項が消える。
Gamow state では実軸上で発散するため、$k$ と輪郭を bound-state 領域から
同じ解析枝で続けるか、Gaussian damping 後に解析接続する。得られる有限値
は Jost 関数の $F'(k_R)$ と同じ留数規格化である。}

\chapter{Stage III：exact WKB と大域接続}

\solproblem{III.1　Riccati 係数}{
対数 WKB ansatz を代入し最初の4係数と枝の parity を求めよ。}{
$[-\hbar^2\partial_x^2+Q(x)]\psi=0$ に
$\psi=\rme^{\hbar^{-1}\int^xS(x',\hbar)\dd x'}$ を代入すると
\begin{equation}
 S^2+\hbar S'=Q,
 \qquad S=\sum_{n\geq0}\hbar^nS_n.
\end{equation}
従って $S_0^2=Q$ と
\begin{equation}
 S_n=-\frac{1}{2S_0}
 \left(S_{n-1}'+\sum_{j=1}^{n-1}S_jS_{n-j}\right)
\end{equation}
を得る。$S_0^\pm=\pm\sqrt Q$ とすれば
\begin{align}
 S_1&=-\frac{Q'}{4Q},\\
 S_2^\pm&=\pm\left(
 \frac{Q''}{8Q^{3/2}}-\frac{5(Q')^2}{32Q^{5/2}}\right),\\
 S_3&=-\frac{Q'''}{16Q^2}
 +\frac{9Q'Q''}{32Q^3}-\frac{15(Q')^3}{64Q^4}.
\end{align}
sheet 交換で $S_{2j}$ は符号を変え、$S_{2j+1}$ は不変。
$\vsub{S}{odd}=(S^+-S^-)/2$ が偶数次数、
$\vsub{S}{even}=(S^++S^-)/2$ が奇数次数を含み、
$\vsub{S}{even}=-(\hbar/2)\partial_x\log\vsub{S}{odd}$ である。}

\solproblem{III.2　最も単純な lateral Borel 和}{
$\sum_{n\geq0}n!\hbar^{n+1}$ の Borel 変換、二つの lateral 和、差を求めよ。}{
Borel 変換は
\begin{equation}
 \widehat F(t)=\sum_{n\geq0}\frac{n!}{n!}t^n=\frac{1}{1-t}.
\end{equation}
正実軸上の $t=1$ に pole があるので
\begin{equation}
 \mathcal S_{0\pm}F(\hbar)
 =\int_{0}^{\infty\rme^{\pm\rmi0}}
 \frac{\rme^{-t/\hbar}}{1-t}\dd t
\end{equation}
と処方する。実部は principal value、虚部は半留数である。上側 minus
下側を本問の規約とすれば
\begin{equation}
 \mathcal S_{0+}F-\mathcal S_{0-}F
 =2\pi\rmi\rme^{-1/\hbar}.
\end{equation}
符号は lateral path の命名を逆にすれば反転するが、ambiguity の大きさと
action $A=1$ は不変である。}

\solproblem{III.3　単純 turning point と Airy 接続}{
単純 turning point を Airy 方程式へ写し、局所接続行列を求めよ。}{
$Q(x_0)=0$、$Q'(x_0)\neq0$ とし
\begin{equation}
 \zeta(x)=\left[\frac32\int_{x_0}^x\sqrt{Q(t)}\dd t\right]^{2/3},
 \qquad \psi=(\zeta/Q)^{1/4}w
\end{equation}
と置く。leading order では
$\hbar^2w_{\zeta\zeta}=\zeta w$、解は
$\operatorname{Ai}(\zeta/\hbar^{2/3})$ と $\operatorname{Bi}$ である。
dominant/subdominant の係数列を $(\vsub{a}{dom},\vsub{a}{sub})^{\mathsf T}$ と
し、Stokes 線を正向きに横切る規約では dominant 係数は不変で
subdominant 係数が $\pm\rmi$ 倍だけ跳ぶ：
\begin{equation}
 \binom{\vsub{a}{dom}}{\vsub{a}{sub}}_{\!\mathrm{after}}
 =\begin{pmatrix}1&0\\ \pm\rmi&1\end{pmatrix}
 \binom{\vsub{a}{dom}}{\vsub{a}{sub}}_{\!\mathrm{before}}.
\end{equation}
符号は $\sqrt Q$ の sheet と横断方向で決まる。ordering を
(sub,dom) に変えれば転置位置へ移る。いずれも determinant 1 で
Wronskian を保存する。}

\solproblem{III.4　Stokes 行列の path ordering}{
三つの triangular Stokes 行列と propagation 行列を掛け、非可換性を示せ。}{
\begin{equation}
 U(a)=\begin{pmatrix}1&a\\0&1\end{pmatrix},\quad
 L(b)=\begin{pmatrix}1&0\\b&1\end{pmatrix},\quad
 P(p)=\begin{pmatrix}p&0\\0&p^{-1}\end{pmatrix}
\end{equation}
とする。各 determinant は1なので、任意の path-ordered product も
determinant 1 で Wronskian を保存する。しかし
\begin{equation}
 U(a)L(b)=\begin{pmatrix}1+ab&a\\b&1\end{pmatrix},\qquad
 L(b)U(a)=\begin{pmatrix}1&a\\b&1+ab\end{pmatrix},
\end{equation}
また $PUP^{-1}=U(p^2a)$ である。従って、例えば
$U(a_3)L(a_2)P(p)U(a_1)$ は横断順序を入れ替えた積と異なる。
determinant だけでは順序誤りを検出できないため、各 crossing の sheet、
向き、basis ordering を ledger に記録する必要がある。}

\solproblem{III.5　barrier resonance と Gamow 幅}{
二 turning point barrier の leading outgoing 条件と幅を導け。}{
classically allowed well の一周期 action を $I(E)$、escape barrier の
Euclidean action を
$K(E)=\int_{x_1}^{x_2}\sqrt{2m(V-E)}\dd x>0$ とする。Airy 接続を一周
掛けると、leading real quantization は
\begin{equation}
 I(E_n)=2\pi\hbar(n+1/2).
\end{equation}
barrier の小さい transmission probability は
$T(E)=\rme^{-2K(E)/\hbar}$。classical period
$\vsub{T}{cl}=\partial I/\partial E$ ごとに一回 escape を試みるなら
decay rate は $T(E_n)/\vsub{T}{cl}$、従って
\begin{equation}
 \Gamma_n=\frac{\hbar}{\vsub{T}{cl}(E_n)}
 \rme^{-2K(E_n)/\hbar}
\end{equation}
（左右二出口なら各寄与を加える）。all-order exact WKB では $K$ を
Borel resummed open Voros period、$I$ を closed quantum period に置き換え、
Stokes multiplier も含む $\Cincoming(E)=0$ を解く。}

\solproblem{III.6　inverted Rosen--Morse の pole}{
hypergeometric 接続から incoming 係数と共鳴・反共鳴枝を求めよ。}{
$V=U_0/\cosh^2(\beta x)$、$E=\hbar^2k^2/(2m)$、
$s(s+1)=-2mU_0/(\beta^2\hbar^2)$ とする。$\xi=\tanh\beta x$ により
associated Legendre 方程式となり、$x\to+\infty$ で outgoing な解を
$x\to-\infty$ へ接続すると、非零位相を除き
\begin{equation}
 \vsub{A}{in}(k)=
 \frac{\Gamma(1-\rmi k/\beta)\Gamma(-\rmi k/\beta)}
 {\Gamma(-\rmi k/\beta-s)
  \Gamma(-\rmi k/\beta+s+1)}.
\end{equation}
Gamma 関数は零を持たないので、分母の pole が $\vsub{A}{in}=0$ を与える。
$\kappa=[8mU_0/(\beta^2\hbar^2)-1]^{1/2}$ とすれば resonance branch は
\begin{equation}
 \vsup{k_n}{R}=\frac{\beta}{2}[\kappa-
 \rmi(2n+1)],\quad
 \vsup{E_n}{R}=\frac{\hbar^2\beta^2}{8m}
 [\kappa-
 \rmi(2n+1)]^2.
\end{equation}
real potential の反共鳴は
$\vsup{k_n}{AR}=-(\vsup{k_n}{R})^*$。どちらも同じ解析式の別 sheet である。}

\solproblem{III.7　単純零点と二重零点の Laurent 展開}{
$S(E)=N(E)/\Cincoming(E)$ を単純零点と二重零点の周りで展開せよ。}{
$x=E-E_*$ とする。単純零点
$\Cincoming=c_1x+c_2x^2/2+\cdots$、
$N=N_0+N_1x+\cdots$ なら
\begin{equation}
 S(E)=\frac{N_0}{c_1}\frac1x
 +\left(\frac{N_1}{c_1}-\frac{N_0c_2}{2c_1^2}\right)+\order(x).
\end{equation}
二重零点
$\Cincoming=c_2x^2/2+c_3x^3/6+\cdots$ なら
\begin{equation}
 S(E)=\frac{2N_0}{c_2}\frac1{x^2}
 +\left(\frac{2N_1}{c_2}-\frac{2N_0c_3}{3c_2^2}\right)\frac1x
 +\order(1).
\end{equation}
Stage VII の Jordan resolvent の $x^{-2}$ 項と $x^{-1}$ 項がこの二係数に
対応する。ただし行列問題では determinant だけでなく左右 null vector
を含む residue を計算する。}

\solproblem{III.8　cubic metastable potential の Stokes graph}{
二枚の $\sqrt Q$ sheet 上で Stokes graph と outgoing path を構成せよ。}{
例 $Q(x,E)=x^2-gx^3-2E$ の三つの turning point $x_j$ を求め、各点から
\begin{equation}
 \im\int_{x_j}^{x}\sqrt{Q(t,E)}\dd t=0
\end{equation}
を満たす三本の Stokes 曲線を数値積分する。branch cut を turning point
間に置き、sheet $+$ では $S_0=+\sqrt Q$、cut 横断後は sheet $-$ と記録する。
origin 側の regular/subdominant 解から metastable well を通り、barrier
turning points の Airy 行列、cut 交換、外側の Stokes jump を順に掛け、
外側で incoming 係数を0にする path を選ぶ。正解の検査は
(i) 各 crossing の順序が連続、(ii) branch 交換が偶数回なら元 sheet、
(iii) 総行列 determinant が1、(iv) contour の小変形で $\Cincoming$ が不変、
の四点である。図の見た目だけで sheet を判定してはならない。}

\solproblem{III.9　analytic coordinate rotation と Voros period}{
turning points、cycle、積分路を同時に回転したとき quantum period が不変であることを示せ。}{
$x=\rme^{\rmi\theta}y$ と同時に potential parameter と turning points を
解析的に移し、cycle $\gamma$ を $\gamma_\theta$ へ写す。Riccati 1-form は
座標変換則により
\begin{equation}
 S_x(x,E;\hbar)\dd x=S_y(y,E;\hbar)\dd y+dd F
\end{equation}
となる。closed cycle では exact differential の積分は0なので
\begin{equation}
 \VorosPeriod_\gamma=oint_\gamma\vsub{S}{odd}\dd x
 =\oint_{\gamma_\theta}\vsub{S}{odd,\theta}\dd y.
\end{equation}
open path でも endpoint 規格化を同時に変換すれば同じである。変形領域を
pole、turning point、branch cut、Borel singular direction が横切ると
homotopy class または lateral sum が変わり、留数あるいは Stokes
automorphism を追加しなければならない。}

\solproblem{III.10　piecewise constant barrier の transfer matrix}{
transfer-matrix pole と WKB 接続条件を比較せよ。}{
区間 $j$ の波を $A_j\rme^{\rmi k_jx}+B_j\rme^{-\rmi k_jx}$ と書く。
幅 $d_j$ の propagation と境界 matching は
\begin{equation}
 P_j=\begin{pmatrix}\rme^{\rmi k_jd_j}&0\\0&\rme^{-\rmi k_jd_j}\end{pmatrix},
\quad
 I_{j,j+1}=\frac12
 \begin{pmatrix}1+k_j/k_{j+1}&1-k_j/k_{j+1}\\
                 1-k_j/k_{j+1}&1+k_j/k_{j+1}\end{pmatrix}.
\end{equation}
$M=I_{N,N+1}P_N\cdots P_1I_{0,1}$ を path order で作る。左右とも
outgoing の条件を課すと、採用した係数順序に応じて $M_{22}(E)=0$
（または等価な minor）が pole 条件である。電流規格化した interface
行列では $\det M=1$ となり、$|\det M-1|$ を丸め誤差 monitor にする。
smooth barrier の区間数を増やすと、$\log P_j$ の和は WKB propagation
action へ、interface 行列は turning-point connection へ収束する。
pole の比較は同じ $k_j$ branch と outgoing convention で行う。}

\chapter{Stage IV：pole の諸実現と continuum response}

\solproblem{IV.1　flat continuum の Feshbach 自己 energy}{
rank-one $P$ state と有限閾値を持つ flat continuum の自己 energy を求めよ。}{
$H_{PP}=E_d$、$\rho|g|^2=c$ を $\vsub{E}{th}<\epsilon<E_c$ で一定とする。
本書の規約では
\begin{equation}
 \Sigma(z)=c\int_{\vsub{E}{th}}^{E_c}\frac{\dd\epsilon}{z-\epsilon}
 =c\log\frac{z-\vsub{E}{th}}{z-E_c}.
\end{equation}
物理上縁 $z=E+\rmi0$、$\vsub{E}{th}<E<E_c$ では
\begin{equation}
 \Sigma(E+\rmi0)=c\,
 \operatorname{PV}\log\left|\frac{E-\vsub{E}{th}}{E-E_c}\right|
 -\rmi\pi c.
\end{equation}
実部が level shift、$-2\im\Sigma=2\pi c$ が幅である。共鳴 pole は
$z-E_d-\vsup{\Sigma}{II}(z)=0$ を second sheet で解く。cutoff を外すなら実部
の発散を $E_d$ に吸収する renormalization 条件が必要である。}

\solproblem{IV.2　実 dilation から complex scaling へ}{
dilation の $x,p,T$ への作用を導き、解析接続せよ。}{
一次元で $(U(s)\psi)(x)=\rme^{s/2}\psi(\rme^sx)$、$s\in\real$ とすると
$U(s)$ は unitary で
\begin{align}
 U(s)xU(s)^{-1}&=\rme^sx,\\
 U(s)pU(s)^{-1}&=\rme^{-s}p,\\
 U(s)TU(s)^{-1}&=\rme^{-2s}T.
\end{align}
$s=\rmi\theta$ へ解析接続すれば
\begin{equation}
 H_\theta=U(\rmi\theta)HU(\rmi\theta)^{-1}
 =\rme^{-2\rmi\theta}T+V(\rme^{\rmi\theta}x).
\end{equation}
これは一般には unitary 同値でなく非自己共役である。$V$ の dilation
analyticity wedge の内側でのみこの式を作用素族として使える。}

\solproblem{IV.3　outgoing wave の square integrability と cut の回転}{
$\rme^{\rmi kr}$、$\im k<0$ が complex scaling 後に可積分となる角度を求めよ。}{
$k=|k|\rme^{-\rmi\phi}$、$0<\phi<\pi/2$ と書く。$r\mapsto
r\rme^{\rmi\theta}$ 後の絶対値は
\begin{equation}
 \left|\rme^{\rmi k r\rme^{\rmi\theta}}\right|
 =\rme^{-|k|r\sin(\theta-\phi)}.
\end{equation}
従って $\theta>\phi$、かつ potential の解析 wedge 内なら $L^2$。
$E\propto k^2$ なので共鳴の偏角は $\arg E_R=-2\phi$、free continuum は
$\rme^{-2\rmi\theta}[0,\infty)$ へ回転する。pole が cut の上側に露出する
条件は $-2\theta<\arg E_R<0$ である。}

\solproblem{IV.4　free essential spectrum と ABC の compact step}{
運動量空間で free spectrum の回転を示し、relative compact perturbation を説明せよ。}{
Fourier 表示で
\begin{equation}
 (H_{0,\theta}\widehat\psi)(\bm p)
 =\rme^{-2\rmi\theta}\frac{p^2}{2m}\widehat\psi(\bm p).
\end{equation}
従って乗算作用素の essential range から
$\spectrum(H_{0,\theta})=
\rme^{-2\rmi\theta}[0,\infty)$。さらに
$V_\theta(H_{0,\theta}-z)^{-1}$ が compact（または $V_\theta$ が
$H_{0,\theta}$-relative compact）なら Weyl 型安定性により essential
spectrum は変わらない。解析 Fredholm 定理は
$\identity-V_\theta\Resolvent_{0,\theta}(z)$ の逆を meromorphic にし、
回転 cut の外に離散固有値として共鳴を与える。}

\solproblem{IV.5　Berggren--Zel'dovich Gaussian continuation}{
$u(r)\sim\rme^{\rmi kr}$ の Gaussian regularized integral を解析接続せよ。}{
基本積分を $q=k+k'$ として
\begin{equation}
 I_\epsilon(q)=\int_0^\infty
 \rme^{-\epsilon r^2+\rmi qr}\dd r
 =\frac{\sqrt\pi}{2\sqrt\epsilon}
 \rme^{-q^2/(4\epsilon)}
 \operatorname{erfc}\left(-\frac{\rmi q}{2\sqrt\epsilon}\right)
\end{equation}
と評価する。まず $\im q>0$ では $\epsilon\downarrow0$ を通常の意味で取り
$I_0(q)=\rmi/q$。右辺は $q$ の解析関数なので、Stokes sector を指定して
resonance quadrant へ続ける。この continued value を
$\int u_ku_{k'}\dd r$ の外側寄与に用いる。$q=0$ の distributional
continuum 規格化と、$\im q<0$ の Gamow 規格化を同じ極限順序で混同しては
ならない。}

\solproblem{IV.6　Gamow functional の test-space continuity}{
contour deformation 後に Gamow functional が連続となる test space を選べ。}{
通常の Schwartz 空間だけでは指数成長を polynomial seminorm で抑えられない。
そこで Schwartz 空間の核部分空間
\begin{equation}
 \begin{aligned}
 \Phi=\{\phi\in\mathcal S(\real_+):\quad
 p_{a,m,n}(\phi)&=\sup_{r\geq0}
 |\rme^{ar}r^m\partial_r^n\phi(r)|<\infty,\\
 &\hspace{4em}\text{for all }a,m,n\}.
 \end{aligned}
\end{equation}
に countable locally convex topology を入れる。ある admissible ray
$r=\rme^{\rmi\theta}s$ 上で Gamow 解が減衰するなら
\begin{equation}
 \dualpair{u_R^\times}{\phi}
 =\int_0^{\infty\rme^{\rmi\theta}}
 u_R(r)\phi(r)\dd r
\end{equation}
は $|\dualpair{u_R^\times}{\phi}|\leq C p_{a,0,0}(\phi)$ と評価でき連続。
Hilbert norm $\|\phi\|_2$ だけでは局所的に同じ norm を持ちながら tail を
遠方へ移した列を制御できず、この評価が失敗する。}

\solproblem{IV.7　continuum level density と散乱位相}{
resolvent 差の trace から level density を導き、位相との関係を示せ。}{
$\Resolvent=(z-H)^{-1}$ の規約で spectral density は
$-\pi^{-1}\im\Resolvent(E+\rmi0)$。従って
\begin{equation}
 \Delta(E)=-\frac1\pi\im\Tr[
 \Resolvent(E+\rmi0)-\Resolvent_0(E+\rmi0)].
\end{equation}
resolvent 差が trace class なら spectral shift function $\xi(E)$ が存在し
$\Delta(E)=\xi'(E)$。Birman--Krein 公式
$\det\Scattering(E)=\rme^{-2\pi\rmi\xi(E)}$ と
$\det\Scattering=\rme^{2\rmi\vsub{\delta}{tot}}$ の位相規約を揃えると
\begin{equation}
 \Delta(E)=\frac1\pi\frac{\rmd\vsub{\delta}{tot}}{\rmd E}
\end{equation}
（$\xi$ の符号規約を逆にすれば両式とも反転）。積分は端点の全位相変化、
すなわち束縛状態数と threshold correction を含む。}

\solproblem{IV.8　有限基底の複素固有値の分類}{
三つの scaling angle の固有値から bound、resonance、rotated continuum を色なしで分類せよ。}{
各 $\theta_j$ で generalized eigenproblem
$H_{\theta_j}c=E N c$ を同じ精度で解く。固有値を nearest-neighbor では
なく左右固有vector overlap と連続 assignment で対応付ける。
\begin{enumerate}[label=(\alph*)]
 \item $E$ が実軸上で $\theta$ と基底に不変、かつ projector norm が安定：bound state。
 \item $\im E<0$ で $|\rmd E/\rmd\theta|$ が誤差以下、左右 residue も安定：resonance。
 \item $\arg E\simeq-2\theta$、角度差に対し
 $E(\theta_2)\simeq\rme^{-2\rmi(\theta_2-\theta_1)}E(\theta_1)$：continuum pseudostate。
\end{enumerate}
距離だけでなく $\theta$ trajectory、residual norm、condition number、basis
拡大時の projector overlap を表にして判定する。}

\solproblem{IV.9　complex-symmetric 行列の c product}{
$c$ product と通常の biorthogonality を比較し、Hermitian conjugation の失敗例を示せ。}{
$H^{\mathsf T}=H$ なら右固有vector $Hr_n=E_nr_n$ に対する左 vector は
$r_n^{\mathsf T}$ であり
\begin{equation}
 (r_m|r_n)_c=r_m^{\mathsf T}r_n,\qquad
 (E_n-E_m)r_m^{\mathsf T}r_n=0.
\end{equation}
例として
$H=\begin{pmatrix}1&\rmi\\\rmi&-1\end{pmatrix}$ は $H^2=0$ の二次 EP で、
零固有vector $r=(-\rmi,1)^{\mathsf T}$ は
$r^{\mathsf T}r=0$ だが $r^\dagger r=2$。Hermitian norm を使うと EP の
self-orthogonality を見失い、resolvent residue の分母
$l^{\mathsf T}H'r$ を $r^\dagger H'r$ に誤置換して pole 強度を誤る。}

\solproblem{IV.10　三つの regularization の同値性}{
Zel'dovich damping、complex rotation、rigged-space pairing が同じ resonance norm を与える理由を説明せよ。}{
まず bound-state 領域 $\im k>0$ で
$I(k)=\int_0^\infty u_k(r)^2\dd r$ を定義する。この領域では
(i) Gaussian factor を入れてから外す、(ii) 積分 ray を
$\rme^{\rmi\theta}\real_+$ へ回す、(iii) analytic test function との pairing
として読む、の三者は Cauchy 変形で同じである。次に integrand、endpoint、
Jost derivative を同じ branch のまま resonance $k_R$ へ解析接続する。
変形 wedge 内に singularity が入らなければ値は一意で
\begin{equation}
 I_R\propto F'(k_R)
\end{equation}
となる。三者は発散実積分を直接等置するのではなく、収束領域で一致する
一つの解析関数の境界値として一致する。}

\solproblem{IV.11　pole residue の微分公式}{
$D(z)=z-E_d-\Sigma(z)$ の単純零点を展開し、$Z_R$ が確率でないことを示せ。}{
$D(z_R)=0$、$D'(z_R)=1-\Sigma'(z_R)\neq0$ だから
\begin{equation}
 \frac1{D(z)}=\frac{Z_R}{z-z_R}+\order(1),
 \qquad Z_R=[1-\Sigma'(z_R)]^{-1}.
\end{equation}
continued sheet では $\Sigma'$ は複素数である。局所例
$\Sigma(z)=(1+\rmi)z+b$ なら $Z_R=1/[-\rmi]=\rmi$。従って $Z_R$ は
一般に実数でも $[0,1]$ 内でもなく、pole residue または biorthogonal
wave-function renormalization であって Born probability ではない。}

\solproblem{IV.12　Breit--Wigner 位相と time delay}{
unitary pole factor から Lorentzian 位相微分と peak time delay を求めよ。}{
$z_R=E_R-\rmi\Gamma/2$ とし
\begin{equation}
 \vsub{S}{pole}(E)=\frac{E-z_R^*}{E-z_R}=\rme^{2\rmi\delta_R(E)}.
\end{equation}
logarithmを微分すると
\begin{equation}
 \frac{\rmd\delta_R}{\rmd E}
 =\frac{\Gamma/2}{(E-E_R)^2+(\Gamma/2)^2}.
\end{equation}
Wigner--Smith delay は $\tau=2\hbar\rmd\delta/\rmd E$ なので
\begin{equation}
 \tau_R(E)=\frac{\hbar\Gamma}{(E-E_R)^2+(\Gamma/2)^2},
 \qquad \tau_R(E_R)=\frac{4\hbar}{\Gamma}.
\end{equation}
background 位相が急変すれば観測 delay はこの単純項との和になる。}

\solproblem{IV.13　Fano line shape}{
定数 coherent background と単純 pole を加え、maximum と zero を求めよ。}{
位相と全体規格化を吸収すると振幅は
$A(\epsilon)=A_0(q+\epsilon)/(\epsilon+\rmi)$、
$\epsilon=2(E-E_R)/\Gamma$ と書ける。従って
\begin{equation}
 \frac{|A|^2}{|A_0|^2}=\frac{|q+\epsilon|^2}{1+\epsilon^2}.
\end{equation}
real $q$ なら zero は $\epsilon=-q$、有限 maximum は
$\epsilon=1/q$ で値 $1+q^2$。$q=q_R+\rmi q_I$ なら分子は
$(\epsilon+q_R)^2+q_I^2$ なので厳密な zero は消え、
$\epsilon=-q_R$ で $q_I^2/(1+q_R^2)$ が残る。}

\solproblem{IV.14　threshold 近傍で異なる幅の定義}{
$\Gamma_l(E)=\gamma k^{2l+1}$ のとき pole、peak、half width を比較せよ。}{
一 channel model を
\begin{equation}
 D(E)=E-E_0-\Delta(E)+\frac{\rmi}{2}\gamma[2\mu E]^{l+1/2}
\end{equation}
とする。pole は指定 sheet で $D(z_p)=0$ を解く。実軸 peak は
$\partial_E|D(E)|^2=2\re[D'(E)D(E)^*]=0$、half-maximum は
$|D(E_\pm)|^2=2|D(\vsub{E}{peak})|^2$ で定まる。$E_R$ が threshold から十分
遠く $\Gamma(E)$ がほぼ一定なら三定義は一致する。threshold では
$k^{2l+1}$ の branch と $\Gamma'(E)$ が大きく、特に $l=0$ で
$\Gamma\propto\sqrt E$ だから peak は pole の $\re z_p$ からずれ、下側
half point が $E=0$ に衝突して非対称になる。従って $-2\im z_p$、FWHM、
peak position を別々に報告する必要がある。}

\chapter{Stage V：radial exact WKB と monodromy}

\solproblem{V.1　d 次元 radial 方程式}{
第一微分を除き、effective inverse-square 係数を求めよ。}{
$d$ 次元で
\begin{equation}
 \nabla^2=\partial_r^2+\frac{d-1}{r}\partial_r
 -\frac{L^2}{r^2},\qquad L^2Y_l=l(l+d-2)Y_l.
\end{equation}
$\Psi=R_l(r)Y_l$、$u_l=r^{(d-1)/2}R_l$ と置くと
\begin{equation}
 -\frac{\hbar^2}{2m}u_l''+\left[V(r)+
 \frac{\hbar^2}{2mr^2}\left\{l(l+d-2)+\frac{(d-1)(d-3)}4\right\}\right]u_l=Eu_l.
\end{equation}
括弧内は
$[l+(d-2)/2]^2-1/4$。この $-1/4$ と exponential map の Schwarzian が
組み合わさって Langer の完全平方を作る。}

\solproblem{V.2　origin の Frobenius 分類}{
$g/r^2$ に対する二指数と square integrability、追加境界条件を論ぜよ。}{
reduced radial Hilbert space $L^2(\real_+,\dd r)$ で
$-u''+g u/r^2=0$ とすると
\begin{equation}
 \nu_\pm=\frac{1\pm\sqrt{1+4g}}2,\qquad u_\pm\sim r^{\nu_\pm}.
\end{equation}
$\nu>-1/2$ が $L^2$ 条件である。$g\geq3/4$ では $u_+$ のみ可積分で
limit point、境界条件は不要。$-1/4<g<3/4$ では二解とも可積分で limit
circle、self-adjoint extension を選ぶ一つの境界条件が要る。$g=-1/4$
では $r^{1/2}$ と $r^{1/2}\log r$、$g<-1/4$ では fall to the centre が
生じ、下方有界な Hamiltonian のため追加の物理的 renormalization が必要。}

\solproblem{V.3　Langer map の導出}{
$r=\rme^x$ と wave-function rescaling を行い $(l+1/2)^2$ を導け。}{
三次元 reduced radial 方程式を
\begin{equation}
 \left[-\hbar^2\partial_r^2+2m(V-E)+\frac{\hbar^2l(l+1)}{r^2}\right]u=0
\end{equation}
と書く。$r=\rme^x$ では
$\partial_r^2=\rme^{-2x}(\partial_x^2-\partial_x)$。さらに
$u(\rme^x)=\rme^{x/2}\phi(x)$ とすれば
$(\partial_x^2-\partial_x)u=\rme^{x/2}(\phi''-\phi/4)$。従って
\begin{equation}
 \left[-\hbar^2\partial_x^2
 +2m\rme^{2x}(V(\rme^x)-E)
 +\hbar^2(l+1/2)^2\right]\phi=0.
\end{equation}
$1/4$ は Jacobian を unitary に実装した結果であり、経験的置換ではない。}

\solproblem{V.4　三次元 harmonic oscillator の二つの量子化}{
closed exact-WKB cycle と open path＋origin monodromy で spectrum を導け。}{
$V=m\omega^2r^2/2$、Langer angular momentum
$L=\hbar(l+1/2)$ とする。radial momentum
\begin{equation}
 p_r(r)=\sqrt{2mE-m^2\omega^2r^2-L^2/r^2}
\end{equation}
の二正 turning point 間の積分は留数計算で
\begin{equation}
 \int_{r_-}^{r_+}p_r\dd r
 =\frac{\pi E}{2\omega}-\frac{\pi L}{2}.
\end{equation}
二 turning point の Maslov 位相を入れ
$\int_{r_-}^{r_+}p_r\dd r=\pi\hbar(n_r+1/2)$ とすれば
\begin{equation}
 E_{n_rl}=\hbar\omega(2n_r+l+3/2).
\end{equation}
closed cycle では左辺を2倍した量を用いる。open path 法では外側
subdominant 解を origin へ接続し、regular Frobenius 指数 $l+1$ の
monodromy $\rme^{2\pi\rmi(l+1)}$ を課す。同じ Langer residue が
turning-point 位相と origin 位相を分けて上式を再現する。}

\solproblem{V.5　Coulomb bound spectrum}{
radial action から Coulomb spectrum を求め、角運動量位相を最後まで分けよ。}{
$V=-\kappa/r$、$E<0$、$L=\hbar(l+1/2)$ とする。
\begin{equation}
 p_r=\sqrt{2m(E+\kappa/r)-L^2/r^2}.
\end{equation}
二 turning point 間の action は $r=0,\infty$ の留数から
\begin{equation}
 \int_{r_-}^{r_+}p_r\dd r
 =\pi\left(\frac{m\kappa}{\sqrt{-2mE}}-L\right).
\end{equation}
turning-point 位相だけを右辺
$\pi\hbar(n_r+1/2)$ に置き、最後に $L=\hbar(l+1/2)$ を入れると
\begin{equation}
 \frac{m\kappa}{\sqrt{-2mE}}=\hbar(n_r+l+1),\qquad
 E=-\frac{m\kappa^2}{2\hbar^2(n_r+l+1)^2}.
\end{equation}
角運動量の $1/2$ と二 turning point の $1/2$ が主量子数の整数を作る。}

\solproblem{V.6　small-circle monodromy の regulator independence}{
$r=\epsilon\rme^{\rmi\phi}$ の局所 monodromy と $\epsilon$ 非依存条件を求めよ。}{
Frobenius 解 $u\sim r^\nu$ は一周で
$u\mapsto\rme^{2\pi\rmi\nu}u$。二解の local monodromy 行列の固有値は
$\rme^{2\pi\rmi\nu_\pm}$ である。WKB では small circle 上の
$\oint S\dd r$ と connection matrices の積が同じ量を与える。切断計算で
bare monodromy $\vsup{\mathcal M_0}{bare}(\epsilon)$ が $\epsilon$ に依存する
なら
\begin{equation}
 \vsup{\mathcal M_0}{phys}
 =\frac{\vsup{\mathcal M_0}{bare}(\epsilon)}{\vsub{Z}{mon}(\epsilon)},
 \qquad
 \epsilon\frac{\rmd}{\rmd\epsilon}
 \log\vsup{\mathcal M_0}{phys}=0
\end{equation}
を renormalization 条件とする。従って
$\epsilon\rmd\log \vsub{Z}{mon}/\rmd\epsilon=
\epsilon\rmd\log\vsup{\mathcal M_0}{bare}/\rmd\epsilon$。物理入力は
選んだ self-adjoint domain または regular Frobenius branch である。}

\solproblem{V.7　radial oscillator の quartic 摂動と最適化}{
weak quartic perturbation を通常摂動論と optimized monodromy で比較せよ。}{
$H=H_\omega+\lambda r^4$、oscillator length
$b=(\hbar/m\omega)^{1/2}$、$N=2n_r+l$ とする。一次摂動は
\begin{equation}
 E^{(1)}_{n_rl}=\lambda b^4
 \left[(N+3/2)(N+5/2)+2n_r(n_r+l+1/2)\right].
\end{equation}
これは無次元結合 $g=\lambda\hbar/(m^2\omega^3)\ll1$ で有効。
optimized 法では任意の $\Omega$ を導入して
\begin{equation}
 H=H_\Omega+\delta\left[\frac{m}{2}(\omega^2-\Omega^2)r^2+\lambda r^4\right],
\end{equation}
$\delta$ で quantum period と origin monodromy を同じ次数まで展開し、
$\delta=1$ 後に $\partial E_N/\partial\Omega=0$（PMS）または最後の係数0
（FAC）で $\Omega_N$ を選ぶ。$g\ll1$ では $\Omega_N\to\omega$ で通常摂動と
一致し、$g\gtrsim1$ では $\Omega\sim(\lambda\hbar/m^2)^{1/3}$ が
strong-coupling scale を取り込み、固定 $\omega$ 展開より有効である。}

\solproblem{V.8　Cornell potential の turning points}{
$V(r)=-\kappa/r+\sigma r$ の turning points を求め、contour の homotopy を検査せよ。}{
Langer 項を含む $p_r^2=2m[E+\kappa/r-\sigma r]-L^2/r^2$ の零点は
\begin{equation}
 -2m\sigma r^3+2mEr^2+2m\kappa r-L^2=0
\end{equation}
の三根である。係数を scale して companion matrix の固有値として求め、
各 parameter step で最小距離 matching だけでなく root の analytic
continuation により label を保つ。物理 cycle は二正実根を囲み、第三根と
$r=0$ pole を横切らない contour とする。
\begin{checkbox}
parameter を小刻みに変え、(i) root の置換を記録、(ii) contour を二通りに
離散化、(iii) $\oint p_r\dd r$ の差が積分誤差以下、(iv) contour と pole の
最小距離を記録する。差が $2\pi\rmi$ times residue だけ跳べば singularity を
横切ったので homotopy class が変わっている。
\end{checkbox}}

\solproblem{V.9　Yukawa-screened Coulomb の threshold}{
$V(r)=-\kappa\rme^{-\mu r}/r$ で pole が cut に近づく点を追跡せよ。}{
turning points は
\begin{equation}
 2m\left[E+\frac{\kappa\rme^{-\mu r}}r\right]-\frac{L^2}{r^2}=0
\end{equation}
の複素零点で、transcendental なので Newton 法と argument principle を併用する。
$\mu=0$ の Coulomb roots から continuation し、同じ sheet の quantum
period 条件を解く。screening $\mu$ を増やすと最外 bound pole
$k=\rmi\kappa_b$ が $\kappa_b\downarrow0$ へ動く。臨界値 $\mu_c$ は
$E=0$ の zero-energy regular 解が infinity で定数項を失う条件、同値に
scattering length $a_0^{-1}=0$ で決まる。$\mu>\mu_c$ では pole は
negative imaginary $k$ axis の second sheet へ入り virtual state となる。
root 数は argument-principle contour で確認する。}

\solproblem{V.10　座標変換が作る見かけの singularity}{
Jacobian monodromy と物理的 self-adjoint-extension parameter を区別せよ。}{
自由粒子 $-\partial_r^2$ に $r=\rme^x$、$u=\rme^{x/2}\phi$ を施すと、元の
potential が正則でも変換後に $+\hbar^2/4$ と $x\to-\infty$ endpoint が
現れる。同様に $r=y^2$ は $y=0$ に inverse-square 型 Schwarzian 項を作る。
しかし unitary Jacobian と元の domain を同時に写せば、変換後の Frobenius
係数比は一意であり新しい物理 parameter はない。これが Jacobian
monodromy である。これに対し元の differential expression 自体が limit
circle endpoint を持ち、二可積分解の係数比を外部物理から選ぶ場合、その
位相または length scale は genuine self-adjoint-extension parameter である。
判定法は「逆変換で domain が既に一意か」である。}

\chapter{Stage VI：few-body continuum、CDCC、complex scaling}

\solproblem{VI.1　二 cluster projectile の座標と reduced mass}{
fragment--target 座標と二つの reduced mass を導き kinetic energy を分離せよ。}{
$P=c+v$ とし
\begin{equation}
 \bm r=\bm r_v-\bm r_c,\qquad
 \bm R=\frac{m_c\bm r_c+m_v\bm r_v}{m_P}-\bm r_T,\quad m_P=m_c+m_v.
\end{equation}
projectile 重心 $\bm r_P$ から
\begin{equation}
 \bm r_c=\bm r_P-\frac{m_v}{m_P}\bm r,\quad
 \bm r_v=\bm r_P+\frac{m_c}{m_P}\bm r
\end{equation}
なので
$\bm R_c=\bm R-(m_v/m_P)\bm r$、
$\bm R_v=\bm R+(m_c/m_P)\bm r$。全重心を除く正準変換により
\begin{equation}
 T=\frac{\bm P_R^2}{2\mu_R}+\frac{\bm p_r^2}{2\mu_r},\qquad
 \mu_R=\frac{m_Pm_T}{m_P+m_T},\quad
 \mu_r=\frac{m_cm_v}{m_P}.
\end{equation}
cross term は mass-weighted 重心の定義により相殺される。}

\solproblem{VI.2　CDCC coupled radial 方程式}{
三体 reaction Hamiltonian を channel basis へ射影し全項を導け。}{
\begin{equation}
 H=K_R+h_P+U_{cT}(\bm R_c)+U_{vT}(\bm R_v),\quad
 \Psi_N^{JM}=\frac1R\sum_c\chi_c^J(R)\mathcal Y_c^{JM}
\end{equation}
とする。$h_P\Phi_i=\epsilon_i\Phi_i$ と angular orthonormality を用いる。
$K_R$ が $R^{-1}\chi Y_L$ に作用すると
\begin{equation}
 K_R\frac{\chi}{R}Y_L=\frac1R\left[-\frac{\hbar^2}{2\mu_R}
 \left(\frac{\rmd^2}{\rmd R^2}-\frac{L(L+1)}{R^2}\right)\chi\right]Y_L.
\end{equation}
$\Bra{\mathcal Y_c}(H-\vsub{E}{tot})\Ket{\Psi_N}=0$ を取れば
\begin{align}
 &\left[-\frac{\hbar^2}{2\mu_R}
 \left(\frac{\rmd^2}{\rmd R^2}-\frac{L_c(L_c+1)}{R^2}\right)
 +\epsilon_c-\vsub{E}{tot}\right]\chi_c^J(R)\\
 &\quad+\sum_{c'}U_{cc'}^J(R)\chi_{c'}^J(R)=0,
\end{align}
\begin{equation}
 U_{cc'}^J(R)=\Bra{\mathcal Y_c^{JM}}
 [U_{cT}(\bm R_c)+U_{vT}(\bm R_v)]\Ket{\mathcal Y_{c'}^{JM}}_{\xi,\hat R}.
\end{equation}
diagonal は channel potential、off-diagonal は excitation、breakup、
continuum--continuum coupling を全て含む。}

\solproblem{VI.3　momentum bin と強収束}{
bin state を規格化し、step projection の強収束と非 norm 収束を示せ。}{
$\braket{\Phi_{k\alpha}|\Phi_{k'\alpha'}}=
\delta_{\alpha\alpha'}\delta(k-k')$ とする。区間 $I_n$ で
\begin{equation}
 \ket{\vsup{\Phi_{n\alpha}}{bin}}=
 \frac1{\sqrt{N_{n\alpha}}}\int_{I_n}f_{n\alpha}(k)
 \ket{\Phi_{k\alpha}}\dd k,\quad
 N_{n\alpha}=\int_{I_n}|f_{n\alpha}(k)|^2\dd k
\end{equation}
なら非重複 bin は orthonormal。$P_N$ を区分一定関数への直交射影とすると、
単関数の $L^2$ 稠密性から任意の固定 wave packet $g$ に対して
$\|P_Ng-g\|_2\to0$。しかし各有限分割に対し、各 bin 内で平均0となる
規格化関数 $g_N$ を選べば $P_Ng_N=0$、従って
$\|\identity-P_N\|=1$。これは strong だが operator-norm ではない収束である。}

\solproblem{VI.4　geometric Gaussian basis の pseudostate}{
短距離 potential を短い Gaussian basis で対角化し、range 依存性を分類せよ。}{
再現可能な例として $\hbar^2/(2\mu)=1$、
$V(x)=-V_0\rme^{-x^2/a^2}$、even basis
$\phi_n(x)=\rme^{-\alpha_nx^2}$、
$\alpha_n=b_n^{-2}$、$b_n=b_1q^{n-1}$ を取る。行列は解析的に
\begin{align}
 N_{ij}&=\sqrt{\frac{\pi}{\alpha_i+\alpha_j}},\\
 T_{ij}&=\frac{2\alpha_i\alpha_j\sqrt\pi}
 {(\alpha_i+\alpha_j)^{3/2}},\\
 V_{ij}&=-V_0\sqrt{\frac{\pi}
 {\alpha_i+\alpha_j+a^{-2}}}.
\end{align}
$Hc=ENc$ を Cholesky orthogonalization 後に解き、$b_N$ を1.3倍ずつ
増やす。負固有値と内部確率
$\int_{|x|<3a}|\psi|^2\dd x$ が安定するものが bound state。正固有値で
内部確率と phase-sensitive observable が plateau を作るものが
resonance-like stabilized pseudostate。$E\propto b_N^{-2}$ で流れ内部確率が
減るものが nonresonant box pseudostate である。$N$ の condition number と
residual $\|Hc-ENc\|$ も同時に報告する。}

\solproblem{VI.5　open/closed Coulomb channel の境界条件}{
Coulomb open channel と Whittaker closed channel を導き flux factor を検査せよ。}{
$E_c=\vsub{E}{tot}-\epsilon_c$。$E_c>0$ なら
$K_c=\sqrt{2\mu_RE_c}/\hbar$、漸近 radial 方程式は Coulomb 方程式なので
\begin{equation}
 \chi_c\to\frac{\rmi}{2}\left[
 H_{L_c}^{(-)}(\eta_c,K_cR)\delta_{cc_0}
 -\sqrt{\frac{K_{c_0}}{K_c}}S^J_{cc_0}
 H_{L_c}^{(+)}(\eta_c,K_cR)\right].
\end{equation}
$H^{(\pm)}$ の Wronskian から flux は $\pm\hbar K_c/\mu_R$ times 振幅
二乗。従って square-root factor 後の各 outgoing flux は
$K_{c_0}|S_{cc_0}|^2$ に比例し、$\sum_c|S_{cc_0}|^2=1$ が得られる。
$E_c<0$ では $K_c=\rmi\kappa_c$ とし、成長解を捨て
$\chi_c\propto W_{-\eta_c,L_c+1/2}(2\kappa_cR)$ を選ぶ。}

\solproblem{VI.6　closed channel の消去}{
二 channel model で closed component を消去し、threshold 上下の polarization potential を求めよ。}{
\begin{align}
 (E-H_{11})\chi_1&=V_{12}\chi_2,\\
 (E-H_{22})\chi_2&=V_{21}\chi_1
\end{align}
から
\begin{equation}
 [E-H_{11}-\vsub{U}{pol}(E)]\chi_1=0,\qquad
 \vsub{U}{pol}=V_{12}(E-H_{22})^{-1}V_{21}.
\end{equation}
spectral resolution $H_{22}\ket{n}=\epsilon_n\ket{n}$ では
$\vsub{U}{pol}=\sum_nV_{12}\ket{n}\bra{n}V_{21}/(E-\epsilon_n)$。全て
$\epsilon_n>E$ の closed 領域では Hermitian coupling に対し実数で、通常
attractive level shift を与える。continuum threshold 上で
\begin{equation}
 (E-H_{22}+\rmi0)^{-1}=\operatorname{PV}(E-H_{22})^{-1}
 -\rmi\pi\delta(E-H_{22}),
\end{equation}
従って $\im \vsub{U}{pol}=-\pi V_{12}\delta(E-H_{22})V_{21}\leq0$ が open
channel への loss width を表す。}

\solproblem{VI.7　complex-scaled smoothing}{
complex-scaled resolution から breakup response を導き、$\theta$ 依存の相殺を示せ。}{
reaction source $\ket{F}=\sum_nT_n\ket{\Phi_n}$ に対し
\begin{equation}
 S(E)=-\frac1\pi\im\bra{F}(E-h_P+\rmi0)^{-1}\ket{F}.
\end{equation}
analytic test space 上で
$(E-h_P)^{-1}=U^{-1}(\theta)(E-h_P^\theta)^{-1}U(\theta)$。extended
resolution を有限基底で近似すると
\begin{equation}
 S(E)\simeq-\frac1\pi\im\sum_\nu
 \frac{\bra{F}U^{-1}\ket{\Phi_\nu^\theta}
       \bra{\widetilde\Phi_\nu^\theta}U\ket{F}}
      {E-E_\nu^\theta}.
\end{equation}
完全基底なら $\sum_\nu\ket{\Phi_\nu^\theta}
\bra{\widetilde\Phi_\nu^\theta}=\identity$ で、隣接する $U^{-1}U$ が相殺し
元の resolvent matrix element に戻る。有限基底の残留 $\theta$ 依存は
pole と rotated-cut pseudostates の不完全な相殺を測る収束誤差である。}

\solproblem{VI.8　alpha+p+n の Jacobi sets と isospin exchange}{
三 Jacobi set を作り、$p\leftrightarrow n$ 対称性とその破れを示せ。}{
粒子を $(p,n,\alpha)=(1,2,3)$ とし、循環置換 $(i,j,k)$ ごとに
\begin{equation}
 \bm x_i=\bm r_j-\bm r_k,\qquad
 \bm y_i=\bm r_i-\frac{m_j\bm r_j+m_k\bm r_k}{m_j+m_k}
\end{equation}
を取れば、$p+n$、$n+\alpha$、$\alpha+p$ pair を内部に持つ三 set が得られる。
isospin limit $m_p=m_n$、$V_{\alpha p}^{N}=V_{\alpha n}^{N}$ では exchange
$P_{pn}$ が $n\alpha$ set と $p\alpha$ set を入れ替え、basis を
$[\Phi_{n\alpha}\pm\Phi_{p\alpha}]/\sqrt2$ に組める。Coulomb
$V_{\alpha p}^{C}$、proton--target Coulomb、電磁 spin-orbit、実際の
$m_p-m_n$ がこの対称性を破る。核力部分と電磁部分を別々に行列化すれば
破れの大きさを直接検査できる。}

\solproblem{VI.9　p+p+7Be の二 rearrangement と double counting}{
同一 proton を反対称化した振幅を作り、inclusive sum の二重計数を検査せよ。}{
labelled amplitude を
$A_1=A[p_1+(p_2{}^7\mathrm{Be})]$、
$A_2=A[p_2+(p_1{}^7\mathrm{Be})]$ とする。proton pair spin $S$ について
\begin{equation}
 A_S(1,2)=\frac1{\sqrt2}
 [A_1(1,2)-(-1)^{S+1}A_2(1,2)]
\end{equation}
と組み、isospin $T=1$ の対称性を含め全 wave function を反対称にする。
物理 phase space を labelled variables で全領域積分するなら $1/2!$ を掛ける。
数値試験は (i) $|A_S(2,1)|^2=|A_S(1,2)|^2$、(ii) 全領域の $1/2!$ 付き
積分が $E_1\geq E_2$ の半領域積分と一致、(iii) coupling を切った極限で
incoherent labelled sum のちょうど半分、を確認する。interference 項を
落とした上で $1/2$ だけ掛けるのは反対称化ではない。}

\solproblem{VI.10　zero recoil での s-wave enhancement}{
spherical Bessel transform を小 recoil momentum で展開し、局所 bin の partial-wave hierarchy を示せ。}{
multipole $l$ の Fourier--Bessel transform は
\begin{equation}
 F_l(q)=\int_0^\infty r^2u_l(r)j_l(qr)\dd r,\qquad
 j_l(z)=\frac{z^l}{(2l+1)!!}[1+\order(z^2)].
\end{equation}
従って $F_0(q)=\int r^2u_0\dd r+\order(q^2)$ だが
$F_1(q)=q\int r^3u_1\dd r/3+\order(q^3)$。$q\simeq0$ の狭い acceptance
では $p$ wave 強度は $q^2$ で抑えられ、small integrated spectroscopic
weight の $s$ wave が支配し得る。一方広い $q$ 範囲では phase space と
大きな $p$ component の積分が勝つ。従って zero-recoil bin の比と全積分
の比は同じ observable ではない。}

\solproblem{VI.11　alpha+n+n の Jacobi 変換と neutron exchange}{
三 Jacobi set、reduced mass、$S=0,1$ の交換符号を求めよ。}{
$m_n=m$、$m_\alpha=M$ とする。$T$ set は
\begin{equation}
 \bm x_T=\bm r_1-\bm r_2,\quad
 \bm y_T=\bm r_\alpha-\frac{\bm r_1+\bm r_2}{2},\quad
 \mu_{xT}=\frac m2,\quad \mu_{yT}=\frac{2mM}{2m+M}.
\end{equation}
二つの $Y$ set は
\begin{align}
 \bm x_1&=\bm r_\alpha-\bm r_1,&
 \bm y_1&=\bm r_2-\frac{M\bm r_\alpha+m\bm r_1}{M+m},\\
 \bm x_2&=\bm r_\alpha-\bm r_2,&
 \bm y_2&=\bm r_1-\frac{M\bm r_\alpha+m\bm r_2}{M+m},
\end{align}
$\mu_x=mM/(m+M)$、$\mu_y=m(m+M)/(2m+M)$。neutron isospin $T=1$ は
対称なので、全体反対称性から spin singlet $S=0$（反対称 spin）には
exchange-even spatial combination、triplet $S=1$ には exchange-odd
combinationを用いる。$T$ set ではそれぞれ $(-1)^{l_x}=+1,-1$ である。}

\solproblem{VI.12　6He response と pole term の符号}{
extended complex-scaled resolution から response を導き、単独 pole term が負になり得る理由を述べよ。}{
$h_6^\theta$ の biorthogonal resolution を挿入して
\begin{align}
 S_F(E)&=-\frac1\pi\im\sum_\nu
 \frac{M_\nu^\theta\widetilde M_\nu^\theta}{E-E_\nu^\theta},\\
 M_\nu^\theta&=\bra{F}U^{-1}\ket{\Phi_\nu^\theta},\\
 \widetilde M_\nu^\theta&=\bra{\widetilde\Phi_\nu^\theta}U\ket{F}.
\end{align}
を得る。isolated pole の積
$M_R^\theta\widetilde M_R^\theta$ は一般に複素で、background との干渉を
含まないため $-\im[M_R\widetilde M_R/(E-z_R)]/\pi$ は局所的に負にもなる。
一方完全な response は元の自己共役 $h_6$ の spectral measure
$\bra{F}\delta(E-h_6)\ket{F}\geq0$ であり、全 pole と cut の coherent sum
後に非負となる。}

\solproblem{VI.13　Riesz projector の pseudostate 表現}{
complex-scaled pole projector を実 pseudostate basis へ戻し、収束試験を与えよ。}{
単純 pole で
\begin{equation}
 P_R^\theta=\ket{\Phi_R^\theta}\bra{\widetilde\Phi_R^\theta},\qquad
 \mathcal P_R=U^{-1}P_R^\theta U.
\end{equation}
$P_N=\sum_i\ket{\Phi_i}\bra{\Phi_i}$ を両側に挿入すると
\begin{equation}
 (\mathcal P_R^{(N)})_{ij}=
 \bra{\Phi_i}U^{-1}\ket{\Phi_R^\theta}
 \bra{\widetilde\Phi_R^\theta}U\ket{\Phi_j},
\end{equation}
すなわち outer product の coherent rank-one 行列を得る。
\begin{checkbox}
基底と $\theta$ を変え、(i) singular value が一つだけ安定、
(ii) $|\Tr\mathcal P_R^{(N)}-1|$、
(iii) $\| (\mathcal P_R^{(N)})^2-\mathcal P_R^{(N)}\|$、
(iv) 物理 source/detector 間の matrix element、を検査する。energy window
内の対角和は rank と位相を保存しないので代用にならない。
\end{checkbox}}

\solproblem{VI.14　三体 phase-space Jacobian}{
Jacobi momentum normalization を宣言して二重微分 breakup strength の Jacobian を導け。}{
$\braket{\bm k_x,\bm k_y|\bm k_x',\bm k_y'}=
\delta^3(\bm k_x-\bm k_x')\delta^3(\bm k_y-\bm k_y')$ とする。
$\epsilon_x=\hbar^2k_x^2/(2\mu_x)$、
$\epsilon_y=\hbar^2k_y^2/(2\mu_y)$ だから
\begin{equation}
 \dd^3k_x\dd^3k_y
 =\frac{\mu_x\mu_y}{\hbar^4}k_xk_y
 \dd\epsilon_x\dd\epsilon_y\rmd\Omega_x\rmd\Omega_y.
\end{equation}
従って本体の $\mathcal J$ は、状態規格化の $(2\pi)$ factors と incident
flux を別にすれば $\mu_x\mu_yk_xk_y/\hbar^4$。二重微分式を両 energy と
両角度で積分すれば元の
$\int|T(\bm k_x,\bm k_y)|^2\dd^3k_x\dd^3k_y$ が再現される。別 Jacobi set
を混ぜるときは mass-scaled 変換の unit determinant と同一粒子 $1/2!$ を
同じ規約で扱う。}

\solproblem{VI.15　二 channel incoming block の pole residue}{
具体的左右 null vector で matrix inverse の pole を検証し、source/detector overlap を同定せよ。}{
\begin{equation}
 \Cincoming(E)=\begin{pmatrix}E-z_R&g\\0&1\end{pmatrix}
\end{equation}
を取る。$E=z_R$ で右 null vector $u_R=(1,0)^{\mathsf T}$、左 null vector
$v_R^\dagger=(1,-g)$、かつ
$v_R^\dagger\Cincoming'(z_R)u_R=1$。実際
\begin{equation}
 \Cincoming(E)^{-1}=
 \frac{u_Rv_R^\dagger}{E-z_R}+
 \begin{pmatrix}0&0\\0&1\end{pmatrix}.
\end{equation}
一般の場合も pole part は
$u_Rv_R^\dagger/[(E-z_R)v_R^\dagger\Cincoming'(z_R)u_R]$。CDCC が作る
channel source column を $s$、CSLS detector row を $d^\dagger$ とすれば
residue は
\begin{equation}
 \frac{(d^\dagger u_R)(v_R^\dagger s)}
 {v_R^\dagger\Cincoming'(z_R)u_R}.
\end{equation}
前二因子がそれぞれ detector と source の overlap、分母が exact-WKB
connection derivative に対応する。}

\chapter{Stage VII：exceptional point と幾何学的 holonomy}

\solproblem{VII.1　2-by-2 complex-symmetric EP}{
一般 complex-symmetric 行列を対角化し二次 EP 条件を求めよ。}{
\begin{equation}
 H=\begin{pmatrix}a&b\\b&d\end{pmatrix},\qquad
 E_\pm=\frac{a+d}{2}\pm\frac12\sqrt{(a-d)^2+4b^2}.
\end{equation}
従って discriminant
\begin{equation}
 \Delta=(a-d)^2+4b^2=0
\end{equation}
で固有値が合流する。さらに $b=0$ かつ $a=d$ なら $H=a\identity$ で二本の
固有vector が残る ordinary degeneracy である。EP には
$H-\vsub{E}{EP}\identity\neq0$、rank 1 が必要。例えば
$r=(b,E-a)^{\mathsf T}$ の二枝は $\Delta\to0$ で同一直線へ合流する。}

\solproblem{VII.2　Jordan chain と resolvent の二重 pole}{
長さ2の Jordan chain を作り resolvent の $x^{-2},x^{-1}$ 項を求めよ。}{
$(H-z_0)r_0=0$、$(H-z_0)r_1=r_0$ とし、Jordan block 上で
$H=z_0\identity+N$、$N^2=0$。$x=z-z_0$ なら本書の resolvent は
\begin{equation}
 (z-H)^{-1}=(x\identity-N)^{-1}
 =\frac{\identity}{x}+\frac{N}{x^2}.
\end{equation}
dual Jordan basis $l_0,l_1$ を
$\dualpair{l_i}{r_j}=\delta_{ij}$ と取れば
$\identity_J=\ket{r_0}\bra{l_0}+\ket{r_1}\bra{l_1}$、
$N=\ket{r_0}\bra{l_1}$。従って
\begin{equation}
 \Resolvent_J(z)=\frac{\ket{r_0}\bra{l_1}}{x^2}
 +\frac{\ket{r_0}\bra{l_0}+\ket{r_1}\bra{l_1}}{x}.
\end{equation}
$\Cincoming(E)$ の二重零点の逆に現れる二係数と同じ pole order である。}

\solproblem{VII.3　EP の self-orthogonality}{
coalesced 右 vector と左 partner の内積が0になることを示せ。}{
EP で右 Jordan chain を
$(H-E_*)r_0=0$、$(H-E_*)r_1=r_0$、左固有vectorを
$l_0(H-E_*)=0$ とする。第二式を左から掛ければ直ちに
\begin{equation}
 \dualpair{l_0}{r_0}=\dualpair{l_0}{(H-E_*)r_1}=0.
\end{equation}
parameter 微分でも同じ結果が見える。
$(H-E)r=0$ を $\lambda$ で微分して左固有vectorを掛けると
$E'\dualpair l r=\dualpair l{H'r}$。EP 近傍で
$E'\sim(\lambda-\lambda_*)^{-1/2}$ が発散する一方、右辺が有限なら
$\dualpair l r\to0$ が必要である。}

\solproblem{VII.4　branch point の一周と二周}{
square-root EP を周回して固有値・固有vector・gauge を追跡せよ。}{
局所 normal form $H(\lambda)=\begin{pmatrix}0&1\\\lambda&0\end{pmatrix}$ を
用いる。$\lambda=\rho\rme^{\rmi\varphi}$ とすると
$E_\pm=\pm\sqrt\rho\rme^{\rmi\varphi/2}$、
$r_\pm=(1,E_\pm)^{\mathsf T}$。$\varphi:0\to2\pi$ で
$E_+\to E_-$、$r_+\to r_-$、すなわち label が交換する。二周で固有値は
元へ戻る。vector の比較には各 step で biorthogonal overlap を正実にする
parallel-transport gauge を選び、最後だけ初期 gauge と比較する。任意に
各点の第一成分を実にする gauge は途中で不連続となり位相を捏造する。
二周後の符号または位相は、この連続 gauge で初めて幾何学的 holonomy と
して定義される。}

\solproblem{VII.5　biorthogonal Berry connection}{
二準位 model の Berry connection と gauge 変換部分を求めよ。}{
上の model で右 vector $r=(1,E)^{\mathsf T}$、左 vector を
$l^{\mathsf T}=(E,1)/(2E)$ とすれば $l^{\mathsf T}r=1$。従って
\begin{equation}
 \mathcal A=\rmi l^{\mathsf T}\rmd r
 =\frac{\rmi}{2}\frac{\rmd E}{E}
 =\frac{\rmi}{4}\frac{\rmd\lambda}{\lambda}.
\end{equation}
$r\mapsto\rme^{f}r$、$l\mapsto\rme^{-f}l$ では
$\mathcal A\mapsto\mathcal A+
\rmi\rmd f$。従って局所 connection は gauge
依存だが、閉路積分の exponential は $f$ が一価なら不変。EP の一周では
state 自体が相手の sheet へ移るので、二周または permutation matrix を
含む non-Abelian transport として比較する。}

\solproblem{VII.6　momentum bin と continuum normalization}{
delta-normalized state を幅 $\Delta k$ の bin に置き換え、発散 normalization と有限 holonomy を分離せよ。}{
$\braket{k|k'}=\delta(k-k')$ とし
\begin{equation}
 \ket{n}=\frac1{\sqrt{\Delta k}}
 \int_{k_0-\Delta k/2}^{k_0+\Delta k/2}\ket{k}\dd k.
\end{equation}
なら $\braket{n|n}=1$。一方、中央の generalized state に戻す係数は
$\ket{n}\simeq\sqrt{\Delta k}\ket{k_0}$、従って
$\ket{k_0}\simeq(\Delta k)^{-1/2}\ket{n}$ と発散する。parameter-dependent
bin の Berry connection は
\begin{equation}
 \dualpair{\widetilde n}{\rmd n}
 =-\frac12\rmd\log\Delta k+\vsub{\mathcal A}{shape}+\vsub{\mathcal A}{phys}.
\end{equation}
最初の exact differential は normalization artifact。閉路で bin の幅と
端点を同じ値へ戻し、shape 項を同一規約で差し引けば
$\oint\mathcal \vsub{A}{phys}$ が有限 holonomy として残る。}

\solproblem{VII.7　bin-width independence の RG 方程式}{
renormalized geometric phase の $\Delta k$ 非依存条件と積分定数を求めよ。}{
bare phase を $\gamma_B(\Delta k)$、counterterm を $C(\Delta k,\mu)$ とし
\begin{equation}
 \gamma_R(\mu)=\gamma_B(\Delta k)+C(\Delta k,\mu),\qquad
 \Delta k\frac{\rmd\gamma_R}{\rmd\Delta k}=0
\end{equation}
を課す。もし $\Delta k\rmd\gamma_B/\rmd\Delta k=c+o(1)$ なら
\begin{equation}
 C=-c\log\frac{\Delta k}{\mu}+C_0,\qquad
 \mu\frac{\rmd\gamma_R}{\rmd\mu}=c.
\end{equation}
$C_0$ は局所発散からは決まらない。reference loop の phase を0とする、
detector の asymptotic phase convention を固定する、または一つの実測値へ
match する、といった物理的 convention/data がこれを決める。}

\solproblem{VII.8　EP と continuum threshold の競合}{
EP が threshold へ近づくとき二つの branch をどう扱うべきか説明せよ。}{
threshold $\vsub{E}{th}$ では resolvent 自己 energy が
$\Sigma(E)\sim a+b(E-\vsub{E}{th})^{l+1/2}$ を持つ。一方孤立 EP は
$E_\pm-E_*\sim\pm c\sqrt{\lambda-\lambda_*}$。両者が近いと局所 determinant
は例えば
\begin{equation}
 D(E,\lambda)=A(E-E_*)^2+B(\lambda-\lambda_*)
 +C(E-\vsub{E}{th})^{l+1/2}+\cdots
\end{equation}
となり、単純な二次多項式 normal form は一様でない。$s$ wave なら
$k=\sqrt{E-\vsub{E}{th}}$ で uniformize し、$D(k,\lambda)=0$ と
$\partial_kD=0$ を同じ sheet で解く。EP sheet exchange と $k\mapsto-k$ の
threshold exchange が合成され、周回 permutation は経路がどの branch
point を囲むかに依存する。}

\solproblem{VII.9　condition number と pseudospectrum}{
EP 周回上の eigenvalue condition number を計算し pseudospectral sensitivity と結べ。}{
単純固有値の condition number は
\begin{equation}
 \kappa_n=\frac{\|r_n\|\|l_n\|}{|\dualpair{l_n}{r_n}|}.
\end{equation}
biorthonormal gauge でも vector norm の積として同じ値が残る。二次 EP から
距離 $\rho=|\lambda-\lambda_*|$ では左右 overlap が
$\order(\sqrt\rho)$ なので $\kappa\sim\rho^{-1/2}$。摂動 $\delta H$ に対し
$|\delta E_n|\lesssim\kappa_n\|\delta H\|$。従って周回上で $\kappa$ が大きい
点は $\epsilon$-pseudospectrum の等高線が膨らむ点と一致する。数値的には
左右 residual を一定に保って $\kappa(\varphi)$ と最小 singular value
$s_{\min}(E-H)$ を同時に plot する。}

\solproblem{VII.10　multiple-zero continuation test}{
$\mathcal C=\partial_E\mathcal C=0$ を用いる探索と near crossing との識別法を設計せよ。}{
scalar determinant または適切な incoming minor について
\begin{equation}
 F_1(E,\lambda)=\mathcal C(E,\lambda)=0,\qquad
 F_2(E,\lambda)=\partial_E\mathcal C(E,\lambda)=0
\end{equation}
を complex Newton 法で同時に解く。Jacobian は
$\partial(E,\lambda)(F_1,F_2)$、必要なら complex-step/automatic
differentiationを使う。候補点の検証は
\begin{enumerate}[label=(\alph*)]
 \item operator の geometric multiplicity が1、Jordan residual が小さい；
 \item 小円周回で二 pole が交換し、二周で戻る；
 \item splitting が radius の平方根に比例；
 \item basis・contour・complex-scaling angle を変えて点と Riesz projector が安定；
 \item $\mathcal C_{EE}\neq0$、$\mathcal C_\lambda\neq0$。
\end{enumerate}
near crossing は有限最小 gap を持ち、一周で label 交換せず、Newton 解が
discretization とともに漂う。}

\chapter{Stage VIII：Regge--Wheeler 方程式と black-hole QNM}

\solproblem{VIII.1　Schwarzschild tortoise 座標}{
Schwarzschild metric から $r_*$ と両端の漸近形を導け。}{
$f(r)=1-2M/r$ とし radial null line $0=-f\dd t^2+f^{-1}\dd r^2$ を
$t\mp r_*=\mathrm{const.}$ にするには $\dd r_*/\dd r=f^{-1}$。従って
\begin{equation}
 r_*=r+2M\log\left(\frac r{2M}-1\right)+C.
\end{equation}
$C=0$ とすると horizon $r=2M+\epsilon$ で
\begin{equation}
 r_*=2M\log\frac{\epsilon}{2M}+2M+O(\epsilon)\to-\infty,
\end{equation}
infinity で
$r_*=r+2M\log(r/2M)+O(M^2/r)\to+\infty$。従って外部領域は
$r_*\in\real$ 全体へ写る。}

\solproblem{VIII.2　spherical harmonics の parity と odd metric 成分}{
scalar/vector/tensor harmonic の parity を示し odd sector の自由成分を数えよ。}{
parity map $(\vartheta,\varphi)\mapsto(\pi-\vartheta,\varphi+\pi)$ で
$Y_{lm}\mapsto(-1)^lY_{lm}$。polar vector $D_AY$ は $(-1)^l$、sphere volume
formを掛けた axial vector $X_A=-\varepsilon_A{}^BD_BY$ は追加の負号を得て
$(-1)^{l+1}$。polar/axial tensor もそれぞれ $(-1)^l,(-1)^{l+1}$。
fixed $(l,m)$ の odd metric は
$h_tX_A$、$h_rX_A$、$h_2X_{AB}$ の3 scalar amplitude。odd gauge function
$\xi X_A$ 一つで $h_2=0$ にでき、2 amplitude が残る。Einstein constraint
が一つを nondynamical にし、最終的な propagating master field は1。
$l=0$ に odd harmonic はなく、$l=1$ で $X_{AB}=0$ なので別扱い。}

\solproblem{VIII.3　linearized Einstein 方程式から covariant Regge--Wheeler 方程式へ}{
warped-product identities を用いた projection と elimination を示せ。}{
$\rmd s^2=g_{ab}\dd x^a\dd x^b+r^2\Omega_{AB}\dd\theta^A\dd\theta^B$ とし、
odd perturbation を $h_{aB}=h_aX_B$、$h_{AB}=h_2X_{AB}$ と展開する。
gauge invariant combination は
\begin{equation}
 \widetilde h_a=h_a-\frac12\nabla_ah_2+\frac{r_a}{r}h_2.
\end{equation}
$D^AX_A=0$、$D^2X_A=[1-l(l+1)]X_A$、
$D^BX_{AB}=-(l-1)(l+2)X_A/2$ を linearized Ricci tensor へ代入すると
\begin{align}
 0=P_a={}&-\Box_2\widetilde h_a+\nabla_a\nabla_b\widetilde h^b
 +\frac2r(r_b\nabla_a\widetilde h^b-r_a\nabla_b\widetilde h^b)\\
 &-\frac2{r^2}r_ar_b\widetilde h^b+\frac{l(l+1)}{r^2}\widetilde h_a,\\
 0=P={}&\nabla_a\widetilde h^a.
\end{align}
$\lambda=(l-1)(l+2)$ とし
\begin{equation}
 \vsub{\Psi}{odd}=\frac{2r}{\lambda}\varepsilon^{ab}
 \left(\nabla_a\widetilde h_b-\frac2r r_a\widetilde h_b\right)
\end{equation}
を定義する。$\varepsilon^{ab}\nabla_aP_b$ を取り、背景 identity
$\nabla_a\nabla_br=(M/r^2)g_{ab}$、$r_ar^a=f$、constraint $P=0$ を使うと
\begin{equation}
 \varepsilon^{ab}\nabla_aP_b=-\frac{\lambda}{2r}
 \left[\Box_2-\frac{l(l+1)}{r^2}+\frac{6M}{r^3}\right]\vsub{\Psi}{odd}.
\end{equation}
従って bracket が0。Bianchi identity は残りの projected equation が従属
であることを保証する。これが成分を飛ばさない gauge-invariant 導出である。}

\solproblem{VIII.4　Regge--Wheeler potential と低 multipole}{
potential を導き両端を検査し、$l=0,1$ の例外を説明せよ。}{
$\Box_2=-f^{-1}\partial_t^2+\partial_r(f\partial_r)$ と
$\partial_{r_*}=f\partial_r$ を前問へ代入し
\begin{equation}
 [-\partial_t^2+\partial_{r_*}^2-\vsup{V_l}{RW}(r)]\vsub{\Psi}{odd}=0,
 \quad \vsup{V_l}{RW}=f\left[\frac{l(l+1)}{r^2}-\frac{6M}{r^3}\right]
\end{equation}
を得る。horizon では $f\to0$ なので $V\to0$ exponentially in $r_*$、
infinity では $V=l(l+1)/r^2+O(r^{-3})\to0$。$l=0$ には odd vector harmonic
が存在しない。$l=1$ では $\lambda=0$ で tensor harmonic と radiative
master definition が消え、vacuum static solution は background の angular
momentumを変える infinitesimal Kerr perturbation である。従って wave
potential formula を機械的に $l=0,1$ へ適用しない。}

\solproblem{VIII.5　horizon と infinity の符号}{
$\rme^{-\rmi\omega t}$ 規約で QNM 条件を vanishing connection coefficient として書け。}{
horizon の future-regular ingoing wave は advanced time $v=t+r_*$ の関数
$\rme^{-\rmi\omega v}$ なので radial factor は
$\rme^{-\rmi\omega r_*}$。infinity の outgoing wave は retarded time
$u=t-r_*$ の関数で radial factor $\rme^{+\rmi\omega r_*}$。horizon-normalized
解を $\vsup{f_H}{in}$ とし infinity で
\begin{equation}
 \vsup{f_H}{in}=\vsub{\mathcal C}{out}(\omega)
 \rme^{+\rmi\omega r_*}+\Cincoming(\omega)
 \rme^{-\rmi\omega r_*}
\end{equation}
と展開する。QNM 条件は $\Cincoming(\omega)=0$。$\im\omega<0$ では両端の
QNM radial factor は実 $r_*$ 上で発散するが、これは符号誤りでなく
resonance boundary condition である。}

\solproblem{VIII.6　barrier-top WKB による fundamental QNM}{
Schwarzschild fundamental mode を leading WKB で見積もり $M$ scaling と比較せよ。}{
frequency equation $[-\partial_{r_*}^2+V]\psi=\omega^2\psi$ の barrier
maximum $r_0$ で $V\simeq V_0+V_0''(r_*-r_{*0})^2/2$ とすると
\begin{equation}
 \omega^2\simeq V_0-
\rmi(n+1/2)\sqrt{-2V_0''}.
\end{equation}
gravitational $l=2$ では
$r_0/M=(9+\sqrt{17})/4=3.280776\ldots$、
$M^2V_0=0.1512867$、$M^4V_0''=-0.00991941$。$n=0$ より
\begin{equation}
 M\vsup{\omega_{20}}{WKB}\simeq0.39885-0.08829\rmi.
\end{equation}
高精度値 $M\omega_{20}\simeq0.37367-0.08896\rmi$ と比べ、damping は良く
実部は leading order の誤差を示す。metric を $r=M\bar r$、
$r_*=M\bar r_*$ と scale すれば方程式は $(M\omega)^2$ のみを含むので
$\omega\propto M^{-1}$。}

\solproblem{VIII.7　tortoise coordinate の complex scaling}{
QNM exposure 条件と inverse map の branch による制限を導け。}{
両端で $r_*\mapsto\rme^{\rmi\theta}r_*$ とする。右端 outgoing wave と左端
ingoing wave（$r_*=-s$）はいずれも magnitude
\begin{equation}
 \left|\rme^{\rmi\omega s\rme^{\rmi\theta}}\right|
 =\rme^{-s\im(\omega\rme^{\rmi\theta})}
\end{equation}
を持つので $\im(\omega\rme^{\rmi\theta})>0$、すなわち
$\theta>-\arg\omega$ が exposure 条件。ただし
\begin{equation}
 r=2M\left[1+W\left(\rme^{r_*/(2M)-1}\right)\right]
\end{equation}
で inverse map は Lambert $W$ の branch を含む。global straight rotation
は logarithm/W branch cut、$r=0$ singularity、他 sheet の horizon image を
横切り得る。従って実際には potential が解析な有限内部区間を保ち、両 tail
だけを曲げる exterior/smooth complex scaling を用いる。}

\solproblem{VIII.8　Gaussian basis 行列と conditioning}{
Regge--Wheeler 方程式の overlap/kinetic 行列を作り精度喪失を診断せよ。}{
一例として中心 $x_j$ と range $b_j$ を持つ
$\phi_j(r_*)=(r_*-x_j)^{p_j}\rme^{-\alpha_j(r_*-x_j)^2}$ を使う。
$p=0$、共通中心なら
\begin{equation}
 N_{ij}=\sqrt{\frac\pi{\alpha_i+\alpha_j}},\qquad
 T_{ij}=\frac{2\alpha_i\alpha_j\sqrt\pi}
 {(\alpha_i+\alpha_j)^{3/2}},\qquad
 V_{ij}=\int\phi_i\vsub{V}{RW}\phi_j\dd r_*.
\end{equation}
complex scaling 後は derivative に $\rme^{-2\rmi\theta}$、potential argument
に deformation を入れる。geometric range ratio $q$ が1に近すぎると columns
がほぼ従属、遠すぎると中間 scale が欠落する。polynomial order 増大も
moment 行列の dynamic range を増やす。
\begin{checkbox}
$N$ の smallest eigenvalue と condition number、Cholesky/SVD cutoff、
generalized residual、左右 overlap、working precision を記録する。precision
を倍にして安定桁が増えない eigenvalue は物理 pole と認定しない。
\end{checkbox}}

\solproblem{VIII.9　Schwarzschild--de Sitter の二 horizon と Nariai limit}{
二 horizon を $r_*\to\pm\infty$ へ写し near-Nariai barrier scaling を求めよ。}{
$f=1-2M/r-\Lambda r^2/3$ の正根を $r_b<r_c$ とする。各単純根で
$f\simeq2\kappa_b(r-r_b)$、
$f\simeq2\kappa_c(r_c-r)$、
$\kappa_h=|f'(r_h)|/2$。従って
\begin{align}
 r_*&\sim\frac1{2\kappa_b}\log(r-r_b)\to-\infty,\\
 r_*&\sim-\frac1{2\kappa_c}\log(r_c-r)\to+\infty.
\end{align}
$\delta=\sqrt{1-9\Lambda M^2}\downarrow0$ と置くと
\begin{align}
 r_{b,c}&=3M\left(1\mp\frac{\delta}{\sqrt3}+O(\delta^2)\right),\\
 r_c-r_b&=2\sqrt3M\delta+O(\delta^2),\\
 \kappa_{b,c}&\sim\frac{\delta}{3\sqrt3M}.
\end{align}
coordinate $r$ では static region が縮むが $r_*$ 幅は
$\order(\kappa^{-1})\sim M/\delta$ に広がる。master potential の高さは
$\order(\kappa^2)$、従って near-Nariai QNM frequency も $\order(\kappa)$。}

\solproblem{VIII.10　QNM pole と continuum level density}{
QNM pole contribution と reference-subtracted cut contribution を比較せよ。}{
適切な source $F$ と reference operator に対し
\begin{equation}
 \Delta_F(E)=-\frac1\pi\im\left[
 \bra{F}\Resolvent(E+\rmi0)\ket{F}-
 \bra{F}\vsub{\Resolvent}{ref}(E+\rmi0)\ket{F}\right].
\end{equation}
contour を変形すると単純 QNM $z_R=\omega_R^2$ の項は
\begin{equation}
 \Delta_R(E)=-\frac1\pi\im\frac{Z_F}{E-z_R},\qquad
 Z_F=\frac{\mathcal N_F(z_R)}{\partial_E\Cincoming(z_R)},
\end{equation}
残りが reference-subtracted branch cut と他 poles。$Z_F$ は複素で cut は
energy dependent なので、全 response の extremum 条件
$\partial_E[\Delta_R+\vsub{\Delta}{cut}]=0$ は一般に
$E=\re z_R$ で満たされない。さらに $E=\omega^2$ の非線形写像により
frequency plot の peak も $\re\omega_R$ と一致しない。pole position、residue、
full real-axis peak を別々に報告する。}

\chapter{付録 R：再総和・最適化・instanton}

\solproblem{R.1　quartic integral の厳密式}{
$I(g)=(2\pi)^{-1/2}\int\rme^{-x^2/2-gx^4/4}\dd x$ を Bessel $K$ で表し、$g\downarrow0$ 展開を検査せよ。}{
被積分関数は偶関数なので $y=x^2$ と置くと
\begin{equation}
 I(g)=\frac1{\sqrt{2\pi}}
 \int_0^\infty y^{-1/2}\rme^{-y/2-gy^2/4}\dd y.
\end{equation}
標準積分
\begin{equation}
 \int_0^\infty\rme^{-a x^4-bx^2}\dd x
 =\frac14\sqrt{\frac ba}\,
 \rme^{b^2/(8a)}K_{1/4}\left(\frac{b^2}{8a}\right)
\end{equation}
へ $a=g/4,b=1/2$ を入れ、負半軸も加えると
\begin{equation}
 I(g)=\frac{\rme^{1/(8g)}}{2\sqrt{\pi g}}
 K_{1/4}\left(\frac1{8g}\right).
\end{equation}
$K_\nu(z)\sim\sqrt{\pi/(2z)}\rme^{-z}
[1+(4\nu^2-1)/(8z)+\cdots]$ を代入すると $I(0^+)=1$、さらに
\begin{equation}
 I(g)\sim\sum_{n\geq0}(-1)^n
 \frac{\Gamma(2n+1/2)}{\sqrt\pi,n!}g^n
 =1-\frac34g+\frac{105}{32}g^2-\cdots
\end{equation}
となり Gaussian moment 展開と一致する。}

\solproblem{R.2　coefficient ratio から Borel singularity}{
三つの連続比から nearest Borel singularity と leading $1/n$ correction を求めよ。}{
係数の厳密比は
\begin{equation}
 r_n=\frac{a_{n+1}}{a_n}
 =-\frac{(2n+1/2)(2n+3/2)}{n+1}
 =-4n-\frac{3}{4(n+1)}.
\end{equation}
$a_n\sim C A^{-n}\Gamma(n)[1+c_1/n+\order(n^{-2})]$ と仮定すると
$A=-1/4$。連続三比 $r_{n-1},r_n,r_{n+1}$ を用いて $A,c_1$ を二 parameter
fit または Richardson extrapolation すれば
\begin{equation}
 a_n\sim C(-4)^n\Gamma(n)
 \left[1-\frac{3}{16n}+\order(n^{-2})\right].
\end{equation}
従って Borel singularity は $t=A=-1/4$、leading $1/n$ coefficient は
$c_1=-3/16$。比 $r_n/(-4n)$ 自体では $1/n$ 項が相殺して最初の補正が
$3/(16n^2)$ になる点に注意する。}

\solproblem{R.3　direct Padé と Borel--Padé at g=0.4}{
両手法を比較し Borel poles と積分可能性を判定せよ。}{
$2m+1$ 個の係数から direct $[m/m]$ と、$b_n=a_n/n!$ に対する Borel
$[m/m]$ を作る。後者を
$g^{-1}\int_0^\infty\rme^{-t/g}P_m(t)/Q_m(t)\dd t$ と積分する。
厳密値は $I(0.4)=0.857608585344409\ldots$。double precision で得る代表値は
\begin{center}
\begin{tabular}{@{}rrrr@{}}
\toprule
$m$ & direct Padé & Borel--Padé & Borel--Padé 誤差\\
\midrule
2&0.8690245430&0.8593359439&$1.73\times10^{-3}$\\
3&0.8623081645&0.8578991758&$2.91\times10^{-4}$\\
4&0.8597745764&0.8576650427&$5.65\times10^{-5}$\\
5&0.8586885057&0.8576207823&$1.22\times10^{-5}$\\
6&0.8581795924&0.8576114484&$2.86\times10^{-6}$\\
8&0.8577902863&0.8576087764&$1.91\times10^{-7}$\\
\bottomrule
\end{tabular}
\end{center}
$m=2$ の Borel poles は約 $-0.3112,-1.0121$、$m=4$ では
$-0.2678,-0.3662,-0.7131,-2.9235$。次数とともに負実軸 $(-\infty,-1/4]$
の cut を模倣し、正実軸上には pole がないので処方なしに積分できる。
正実 pole が現れる off-diagonal sequence は Froissart pair か真の obstruction
かを residue と次数安定性で判定し、理由なく避けて積分してはならない。}

\solproblem{R.4　third-order optimized perturbation}{
$I_N(g,\Omega)$ を導き正の PMS roots を追跡せよ。}{
$A=1-\Omega$ とし、$\Omega$ Gaussian moment
$\braket{x^{2m}}_\Omega=(2m-1)!!\Omega^{-m-1/2}$ を使う。$\delta=1$ 後
\begin{align}
 I_0={}&\Omega^{-1/2},\\
 I_1={}&I_0-\frac{A}{2\Omega^{3/2}}-\frac{3g}{4\Omega^{5/2}},\\
 I_2={}&I_1+\frac{3A^2}{8\Omega^{5/2}}
 +\frac{15Ag}{8\Omega^{7/2}}+\frac{105g^2}{32\Omega^{9/2}},\\
 I_3={}&I_2-\frac{5A^3}{16\Omega^{7/2}}
 -\frac{105A^2g}{32\Omega^{9/2}}
 -\frac{945Ag^2}{64\Omega^{11/2}}
 -\frac{3465g^3}{128\Omega^{13/2}}.
\end{align}
$\partial_\Omega I_N=0$ の全正実根を走査すると $N=2$ には正根がなく、
$N=1,3$ には次の連続枝が各一つある。
\begin{center}
\begin{tabular}{@{}rcc@{}}
\toprule
$g$ & $\vsub{\Omega}{PMS}^{(1)}$ & $\vsub{\Omega}{PMS}^{(3)}$\\
\midrule
0.05&1.112372&1.212272\\
0.10&1.207107&1.374451\\
0.20&1.366025&1.631074\\
0.40&1.618034&2.019426\\
0.60&1.822876&2.327015\\
1.00&2.158312&2.823067\\
\bottomrule
\end{tabular}
\end{center}
root tracking は前の $g$ の解を初期値にし、同時に全正根の bracket scan を
行う。偶数次数で根がないことを隠して complex root を実根として選ばない。}

\solproblem{R.5　ODM の独立実装}{
$g=\rho\lambda/(1-\lambda)^2$ の ODM を実装し root continuity の意味を説明せよ。}{
quadratic equation
$g(1-\lambda)^2=\rho\lambda$ のうち $g\to0$ で $\lambda\to0$ となる枝を
選ぶ。weak series を代入すると
\begin{equation}
 c_0=1,\qquad
 c_n(\rho)=\sum_{k=1}^{n}a_k\rho^k
 \binom{n+k-1}{n-k}\quad(n\geq1),
\end{equation}
従って $c_N(\rho_N)=0$ の正根を全て求め
$I_N=\sum_{n=0}^Nc_n(\rho_N)\lambda(g,\rho_N)^n$ を評価する。実装の検算は
(i) $\lambda$ 展開を再度 $g$ へ戻し $a_0,\dots,a_N$ を再現、
(ii) polynomial residual $|c_N(\rho_N)|$、(iii) working precision、で行う。
order $N$ ごとに別の根を最小誤差だけで選ぶと、写像する analytic domain
自体が不連続に変わり sequence $I_N$ の意味がなくなる。根は $N$ と $g$ の
homotopy で同じ枝を追い、消滅時にはその事実を報告する。}

\solproblem{R.6　dispersion relation から large order}{
instanton discontinuity を代入し $A$ の冪を含めて係数漸近を導け。}{
\begin{align}
 a_n&=\frac{(-1)^{n+1}}\pi\int_0^\infty
 \frac{\im F(-s+\rmi0)}{s^{n+1}}\dd s,\\
 \im F(-s+\rmi0)&\sim Cs^{-b}\rme^{-A/s}.
\end{align}
へ代入する。$u=A/s$、$s=A/u$、$\dd s=-A u^{-2}\dd u$ なので
\begin{align}
 a_n&\sim\frac{(-1)^{n+1}C}{\pi}
 \int_0^\infty s^{-n-b-1}\rme^{-A/s}\dd s\\
 &=\frac{(-1)^{n+1}C}{\pi}
 A^{-n-b}\int_0^\infty u^{n+b-1}\rme^{-u}\dd u\\
 &=\frac{(-1)^{n+1}C}{\pi A^{n+b}}\Gamma(n+b).
\end{align}
$A^{-n}$ だけでなく fluctuation prefactor $s^{-b}$ が追加の $A^{-b}$ と
$\Gamma(n+b)$ を作る。}

\solproblem{R.7　Rosen--Morse における instanton action の対応物}{
inverted Rosen--Morse で $A=1/4$ に代わる WKB action と Stokes multiplier を同定せよ。}{
spectral curve
\begin{equation}
 y^2=2m\left[\frac{U_0}{\cosh^2(\beta x)}-E\right]
\end{equation}
の relevant turning points $x_1,x_2$ を、選んだ outgoing sheet 上で結ぶ
open action
\begin{equation}
 A_R(E)=\int_{x_1}^{x_2}y(x,E)\dd x
\end{equation}
（または同じ homology class の半 closed Voros period）が zero-dimensional
例の定数 $1/4$ を置き換える。実際の nonperturbative factor は
$\rme^{-2A_R(E)/\hbar}$ であり、periodicity
$x\mapsto x+\rmi\pi/\beta$ による複数 turning-point image も含めて Borel
singularities を作る。単純 turning point の Stokes multiplier と propagation
$\rme^{\pm A_R/\hbar}$ を path order で掛けた incoming entry が
$\Cincoming(E)$。exact solution ではその積が
\begin{equation}
 \Cincoming(k)\propto
 \frac{\Gamma(1-\rmi k/\beta)\Gamma(-\rmi k/\beta)}
 {\Gamma(-\rmi k/\beta-s)\Gamma(-\rmi k/\beta+s+1)}
\end{equation}
へ resummation される。従って Gamma denominator の pole は、quantum
action と Stokes multiplier を含む exact-WKB incoming coefficient の零点である。}


\backmatter
\chapter*{記号の最終確認}
作用素の等式は指定された定義域上で読む。連続スペクトル上の
$\Resolvent(E\pm\rmi0)$ は境界値であり、$L^2$ 上の通常の逆作用素では
ない。共鳴極はシートと境界条件を指定して初めて意味を持つ。有限基底の
複素固有値は、変形角・基底・切断を変えた安定性試験を通過した場合にのみ
共鳴と同定する。
\end{document}
